% =====================================================================
%  Harald Høffding, HENRI BERGSON'S FILOSOFI. KARAKTERISTIK OG KRITIK.
%  Kjøbenhavn og Kristiania: Gyldendalske Boghandel, Nordisk Forlag, 1914
%  (Græbes Bogtrykkeri).  [VIII] + 70 pp.
%  TRANSCRIPTION (Danish) of the 1914 print, complete.
%
%  GENERATED -- DO NOT EDIT BY HAND.  Rebuild from the repo root:
%    sh ebook-method/bergsons-filosofi/rebuild.sh
%  (draft.py from the e-book + template.tex, check.py --apply with
%  decisions.tsv, check.py --mark --write, draft.py --post with patches.POST).
%
%  First (and only) edition.  Text pp. 1-62 in six chapters with numbered
%  paragraphs (1-23, 26-33: the print's numbering skips 24-25); NOTER
%  (half-title p. 63, p. 64 blank, the 52 endnotes pp. 65-70).
%
%  SOURCE: HathiTrust's PDF (umn.31951002074925j, record 012224438) of
%  Google's scan of the University of Minnesota copy, ~/bibliotek/Høffding,
%  Harald/1914-henri-bergsons-filosofi.pdf (SCANS.tsv, sha256 90dd232d...).
%  PAGE MAP: printed page = PDF page - 13 (p. 1 = PDF 14, p. 70 = PDF 83).
%  Front matter (unnumbered): half-title PDF 8, the list of the author's
%  works PDF 9, title PDF 10, imprint PDF 11, Indhold PDF 12.  Not
%  transcribed: HathiTrust's sheet, the wrapper, library labels and stamps
%  (N. B. M. F.'s Bibliotek No. 5593; University of Minnesota), the
%  publisher's advertisements after p. 70.  Every page turn is marked
%  "% --- p. N ---" + \opage{N} at its first word (check.py --mark: 69 of
%  69 exact); a word divided over a turn ends in %.
%
%  CONVENTIONS: italics \textit, letterspacing \emph (names in the Noter),
%  the first word of each chapter \initial (a large initial, the rest in
%  small capitals); note references \nm{N} (a superscript figure); dashes
%  ---, page ranges --; quotation marks „...“ as printed, including the
%  print's repeated „ at the head of a continued line (p. 61 „geründet“:
%  hence one „ more than “).  Running heads, folios and sheet signatures
%  („H. Høffding: Henri Bergson's Filosofi. 5“) not transcribed.
%
%  METHOD AND ACCURACY (2026-10-27; TRANSCRIPTION-PLAYBOOK §00 point 0, the
%  e-book method).  No one typed the text.  The base is the SAGA e-book
%  (Lindhardt og Ringhof 2023, ISBN 9788726362985; edition check in
%  RESUME-NOTES: the 1914 text, nothing modernised, nothing missing),
%  converted by ebook-method/bergsons-filosofi/draft.py, and collated word
%  by word with punctuation against the scan's own Google text layer
%  (ebook-method/collate/check.py, lines rebuilt from the word boxes,
%  punctuation glued; 21,051 words against 20,962): 261 disagreements.
%  30 settled by rule (22 glyph, 5 word division, 2 moved, 1 speck); 84
%  long dashes the layer drops settled by rule after a BLIND audit of 25 at
%  the image (25 of 25 present); the rest read at the image in BLIND crops
%  (116 crops, two readers who saw neither reading; 4 planted controls, 4
%  read right) and in zooms read in the main session (z1-z7).  The e-book
%  was wrong at: p. 1 „al Mening“, „ligger til Grund“, Udviklingshypotesen,
%  bestrider, Maade; p. 28 „Andet, end“; p. 36 „Det“ (no accent).  The
%  NOTER were read whole and blind at the image and set from that reading
%  (notesread2tex.py), compared word by word with the e-book and the layer
%  (notescmp.py) and settled at zooms: from the reading come the
%  letterspaced names (confirmed by a measure of the white space inside
%  each word box, spacing.py: names 0.24-0.32 against a median of 0.11),
%  Bibliothèque, „p. 50“, „IX. p.“, the edition figure 8 (not 3) after
%  Mémoire in notes 16, 20, 24 (3 in notes 21-23), and the labels' stops;
%  the reader's „lagttagelse“ and an italic „München.“ were his own
%  (patches.NOTES).  Italics: the body is roman but for three spans (two
%  titles, „Vi leve“ p. 58), all read at the image; the stroke-slant
%  measure's flags in the body (words after an opening „) were all upright
%  in blind crops.  No letterspacing in the body (spacing.py: none above
%  the narrow-letter words).  Two pages read whole and blind (pp. 5, 11,
%  seed 20261027): 0 word differences in 570 words.  Unattested-word sweep
%  (160 words found in no other transcription in the collection): 3
%  misprints, below.  The ø and umlaut rules: no sites.
%  Word-accuracy: the whole text diffed against the Google layer
%  (check.py), 261 candidates, 13 e-book errors corrected (7 text, 6 notes), the
%  rest the witness's fault, 0 unresolved; the notes read whole at the image.
%  What this cannot see: a misprint that the e-book and the layer both
%  silently normalised outside the two pages read whole.
%
%  PRINTED AS IS (each marked at its line): p. 7 Fremstiling; p. 27
%  „Matière et Memoire“; p. 28 motrîce; p. 32 Øjeblikkke; p. 33
%  Hjærneprocesserne; Noter: note 14 „c’est-â-dire“, note 18 Tannéry and
%  Brunschwieg (thrice; Brunschvicg), note 44 „(p. 55)“ without its stop.
%  Found while translating (the odd-spot list, read at zooms z8): p. 4
%  dennne; p. 29 „forberedte Hindringer“ for Handlinger (Bergson's „nos
%  actions commencées ou préparées“); note 13 object (objet); note 36
%  „(-rattacher“.
%  Kept as printed, not marked: p. 34 Kant's „bleyernes Waffen“; p. 61
%  Goethe's „geründet“.
% =====================================================================
\documentclass[12pt,a4paper,openany]{book}
\usepackage[utf8]{inputenc}
\usepackage[T1]{fontenc}
\usepackage{libertinus}
\usepackage{libertinust1math}
\usepackage{textalpha}
\usepackage[danish]{babel}
\usepackage[protrusion=true,expansion=true]{microtype}
\usepackage{setspace}
\usepackage[margin=1.3in]{geometry}
\usepackage{fancyhdr}
\usepackage{hyperref}
\hypersetup{pdftitle={Henri Bergson's Filosofi. Karakteristik og Kritik}, pdfauthor={Harald Høffding}, hidelinks}
\setlength{\headheight}{15pt}
\pagestyle{fancy}
\fancyhf{}
\fancyhead[LE,RO]{\thepage}
\fancyhead[C]{\textit{Høffding: Henri Bergson's Filosofi (1914)}}
\setstretch{1.3}
\usepackage{marginnote}
\newcommand{\opage}[1]{\marginnote{\footnotesize\textit{[#1]}}}
% a chapter head as printed: the title in capitals, centred, over a short rule
\newcommand{\chapterhead}[1]{\clearpage\phantomsection\addcontentsline{toc}{chapter}{#1}\vspace*{3cm}\begin{center}{\large #1}\\[6pt]\rule{1.2cm}{0.4pt}\end{center}\medskip}
% the first word of each chapter: a large initial, the rest in small capitals (blind reads O001-O006)
\newcommand{\initial}[2]{{\Large #1}\textsc{#2}}
% the short rule that closes each chapter
\newcommand{\chapterend}{\par\begin{center}\rule{2.5cm}{0.4pt}\end{center}}
% a note reference: a superscript figure; written after a space, which it removes
\newcommand{\nm}[1]{\unskip\textsuperscript{#1}}
% the Noter: their half-title (p. 63) and each note's label
\newcommand{\noterhead}[1]{\phantomsection\addcontentsline{toc}{chapter}{#1}\vspace*{6cm}\begin{center}{\Large #1}\end{center}}
\newcommand{\notelabel}[1]{\noindent #1}
% the Indhold's lines
\newcommand{\ix}[3]{\par\noindent\makebox[2em][l]{#1}#2\dotfill #3}

\begin{document}

\frontmatter
\pagestyle{empty}
% --- half-title (PDF 8) ---
\vspace*{5cm}
\begin{center}
{\Large HENRI BERGSON’S FILOSOFI}
\end{center}
\clearpage

% --- verso of the half-title: the author's works (PDF 9) ---
% read whole at the image, blind (reads/pages-read.txt, PDF 9); set inside a dotted frame in the print, the titles in capitals
\begin{center}
{\large HARALD HØFFDING}
\end{center}
\medskip
\begin{footnotesize}
\setlength{\parindent}{0pt}\setlength{\parskip}{1pt}\setstretch{1.1}
\hangindent=1.5em\hangafter=1 DEN ANTIKE OPFATTELSE AF MENNESKETS VILLIE. 1870.\par
\hangindent=1.5em\hangafter=1 PHILOSOPHIEN I TYSKLAND EFTER HEGEL. 1872.\par
\hangindent=1.5em\hangafter=1 DEN ENGELSKE PHILOSOPHI I VOR TID. 1874.\par
\hangindent=1.5em\hangafter=1 OM GRUNDLAGET FOR DEN HUMANE ETHIK. 1876.\par
\hangindent=1.5em\hangafter=1 PSYKOLOGI i Omrids paa Grundlag af Erfaring (1882). Sjette paa ny gennemsete Udgave. 1911.\par
\hangindent=1.5em\hangafter=1 FORMEL LOGIK. Til Brug ved Forelæsninger. (1884). Sjette Udgave. 1913.\par
\hangindent=1.5em\hangafter=1 ETIK. En Fremstilling af de etiske Principer og deres Anvendelse paa de vigtigste Livsforhold. (1887). Tredie paa ny gennemsete og delvis ændrede Udgave. 1905.\par
\hangindent=1.5em\hangafter=1 PSYKOLOGISKE UNDERSØGELSER. 1889. (Videnskabernes Selskabs Skrifter).\par
\hangindent=1.5em\hangafter=1 ETISKE UNDERSØGELSER. Om Muligheden af en filosofisk Etik. --- Om Velfærdsprincipet. --- Forholdsloven i Etiken. 1891.\par
\hangindent=1.5em\hangafter=1 SØREN KIERKEGAARD SOM FILOSOF. 1892.\par
\hangindent=1.5em\hangafter=1 KONTINUITETEN I KANTS FILOSOFISKE UDVIKLINGSGANG. 1893. (Videnskabernes Selskabs Skrifter).\par
\hangindent=1.5em\hangafter=1 DEN NYERE FILOSOFIS HISTORIE. En Fremstilling af Filosofiens Historie fra Renaissancens Slutning til vore Dage. To Bind. (1894--95). Anden gennemsete Udgave. 1904.\par
% "Slutning" is partly covered by ink strokes in the Minnesota copy (zoom z5); the reading follows the book's known subtitle
\hangindent=1.5em\hangafter=1 JEAN JACQUES ROUSSEAU og hans Filosofi. (1896). Anden Udgave. 1912.\par
\hangindent=1.5em\hangafter=1 KORT OVERSIGT OVER DEN NYERE FILOSOFIS HISTORIE. (1898). Fjerde paany omarbejdede Udgave. 1910.\par
\hangindent=1.5em\hangafter=1 MINDRE ARBEJDER. Første Række. 1899.\par
\hangindent=1.5em\hangafter=1 DET PSYKOLOGISKE GRUNDLAG FOR LOGISKE DOMME. 1899. (Videnskabernes Selskabs Skrifter).\par
\hangindent=1.5em\hangafter=1 VED AARHUNDREDESKIFTET. Tale holdt paa Københavns Universitet. 1901.\par
\hangindent=1.5em\hangafter=1 RELIGIONSFILOSOFI. (1901). Andet Oplag. 1906.\par
\hangindent=1.5em\hangafter=1 FILOSOFISKE PROBLEMER. 1902. (Universitetsprogram).\par
\hangindent=1.5em\hangafter=1 OM NOGLE RELIGIONSFILOSOFISKE ARBEJDER fra den nyeste Tid. 1903. (Universitetsprogram).\par
\hangindent=1.5em\hangafter=1 MODERNE FILOSOFER. 1904.\par
\hangindent=1.5em\hangafter=1 MINDRE ARBEJDER. Anden Række. 1905.\par
\hangindent=1.5em\hangafter=1 DANSKE FILOSOFER. 1909.\par
\hangindent=1.5em\hangafter=1 DEN MENNESKELIGE TANKE, dens Former og dens Opgaver. 1910.\par
\hangindent=1.5em\hangafter=1 RELIGION OG VIDENSKAB. 1910.\par
\hangindent=1.5em\hangafter=1 UDVALGTE STYKKER af dansk filosofisk Literatur. Med Indledninger af \textit{Harald Høffding}. 1910.\par
\hangindent=1.5em\hangafter=1 PERSONLIGHETSPRINCIPEN i FILOSOFIN. Stockholm 1911.\par
\hangindent=1.5em\hangafter=1 MINDRE ARBEJDER. Tredie Række. 1913.\par
\end{footnotesize}
\clearpage

% --- title page (PDF 10; the University of Minnesota's and the N. B. M. F.'s library stamps are not transcribed) ---
\begin{titlepage}
\centering
{\large H\,A\,R\,A\,L\,D\ \ H\,Ø\,F\,F\,D\,I\,N\,G}\par   % letterspaced in the print
\vspace{2cm}
{\LARGE HENRI BERGSON’S\\[4pt] FILOSOFI}\par
\vspace{1cm}
{\large KARAKTERISTIK\\ OG KRITIK}\par
\vfill
{\small KJØBENHAVN OG KRISTIANIA}\par
\medskip
{GYLDENDALSKE BOGHANDEL\\ NORDISK FORLAG}\par
\medskip
1914\par
% (the publisher's device, an oval wreath with the monogram GB, stands above the imprint)
\end{titlepage}

% --- verso of the title (PDF 11) ---
\vspace*{\fill}
\begin{center}
{\footnotesize KØBENHAVN --- GRÆBES BOGTRYKKERI}
\end{center}
\clearpage

% --- the Indhold (PDF 12, one unnumbered page), set as printed ---
\begin{center}{\large INDHOLD}\\[4pt]\rule{1.2cm}{0.4pt}\end{center}
\medskip
\noindent\hfill{\small Side}
\ix{1.}{Filosofiens Opgave}{1}
\ix{2.}{Intuition}{11}
\ix{3.}{Psykologi og Fysiologi}{27}
\ix{4.}{Udviklingsfilosofi}{35}
\ix{5.}{Villiens og Latterens Psykologi}{43}
\ix{6.}{Forsøg paa Metafysik}{52}
\par
\begin{center}\rule{2cm}{0.4pt}\end{center}
\clearpage

\mainmatter
\pagestyle{fancy}
% ===== TEXT (pp. 1--62) and NOTER (pp. 63--70) =====
% --- p. 1 ---
\opage{1}\chapterhead{FILOSOFIENS OPGAVE}

1. \initial{I}{} den sidste Del af det nittende Aarhundrede banede Udviklingstanken sig Vej paa Forskningens forskellige Omraader og afgav nye ledende Synspunkter med Hensyn til den menneskelige Erkendelses Oprindelse og Gyldighed. Naar Faust hos Goethe ikke er tilfreds med den sædvanlige Oversættelse af første Vers i Johannesevangeliet, men istedetfor „Ordet“ sætter „Tanken“ (der Sinn), og saa igen ogsaa bliver betænkelig ved dette Udtryk, fordi det ikke er Tanken eller Meningen, der skaber Alt, og tilsidst oversætter: „I Begyndelsen var Daaden“, --- saa udtrykker det Udviklingslærens Opfattelse. Al Tanke, al Mening, vi kunne finde, udspringer af Virkelighedens Tugt gennem selve de Handlinger, den nøder os til at foretage. Videnskabens Gyldighed faar sin sidste Begrundelse deri, at al vor Videnskab er en Gennemførelse af de Forudsætninger, Menneskene maa gaa ud fra, hvis deres Handlinger skulle føre til de Resultater, de ønske. Den praktiske Nødvendighed ligger til Grund for den teoretiske og godtgør dennes Gyldighed. Den saakaldte Pragmatisme har senere, efterat Drøftelsen af Udviklingshypotesen var sagtnet, paa rent psykologisk og historisk Grundlag ført denne Tankegang videre.

Henri Bergson slutter sig til denne Opfattelse, men bestrider de Slutninger, man drager af den. Den praktiske Nødvendighed af at lægge sine Forestillinger til Rette paa en vis Maade kan i hans Øjne ingenlunde godtgøre Gyldigheden af
% --- p. 2 ---
\opage{2}den Erkendelse, man saaledes naar. Tvertimod. Langtfra at Praxis skulde føre til den sande Teori, fremelsker Livets Krav, Kampen for Existensen i alle dens Former, det stadige Arbejde paa at tvinge de materielle Ting til at tjene Menneskets Behov netop en udvortes, mekanisk Maade at tænke paa, en Udstykning og Spredning af Forestillingerne, der gør det umuligt at vinde Indsigt i Tilværelsens indre Sammenhæng.

Under Kampen med det Materielle smittes Tanken af sin Genstand og bliver da behersket af Rumsbilleder, vænner sig til at opfatte alle Forhold som Rumsforhold og til at antage et ligesaa udvortes Forhold mellem Forestillingerne, som der er mellem de forskellige Steder og Genstande i Rummet. Kun en helt udvortes Sammenknytning bliver da mulig. Selve Sproget, Nødvendigheden af at udtrykke Tankerne i Ord, fører i samme Retning. Hver Ting og hver Egenskab betegnes med sit Ord, og man mener da, at de ere ligesaa forskellige fra hverandre, som de forskellige Ord ere.

Ikke blot Naturvidenskaben, Læren om det Materielle, har ligget under for denne Indflydelse, men ogsaa Filosofien. Helt tilbage i den græske Filosofi sporer Bergson den, og det er hans Overbevisning --- som jeg senere hen nærmere skal drøfte --- at den nyere Filosofi i denne Henseende ikke har frigjort sig fra den antike. De ældste græske Filosofer holdt sig meget nær ved Anskuelsen, saaledes som denne umiddelbart fremstiller Tingene; de antog, at det var et af de for Sansningen tilgængelige Elementer, som ved Omsætning eller Forvandling gik over til de andre Elementer. Men med Parmenides fra Elea (ca. 500 f. Kr.) kom Filosofien ind paa den Bane, den siden ikke har forladt. De eleatiske Filosofers Kritik af Begrebet Forandring og deres Hævden af det Uforanderlige som det, der ene kan være Genstand for sand Viden, førte til at hæve det rene Begreb paa Anskuelsens Bekostning, til at bevæge sig i Abstraktionens, Generalisationens og Slutningernes Verden istedetfor at gaa op i, hvad Sansning og umiddelbar Bevidsthed indeholde. I Platon’s Idelære fastsloges de rene Begreber som Tilværelsens inderste Væsen i Modsætning til den stadige Skiften indenfor Sanseanskuelsernes og
% --- p. 3 ---
\opage{3}Erfaringernes Verden \nm{1}. Og siden er, ifølge Bergson, Filosofien stadigt gaaet i samme Retning.

Heri er der nu, ifølge Bergson, en stor Miskendelse af Filosofiens Betydning, som strækker sig gennem hele dens Historie. Denne Betydning bestaar, hævder han, netop i at frigøre Tænkningen, og Sjælelivet idethele, fra den Forkærlighed for det Uforanderlige, for det, der stadigt gentager sig, for det Faste og Mekaniske, som den praktiske Behandling af de materielle Ting fører til at fremkalde. Videnskaben og Filosofien have hidtil staaet i praktiske Formaals Tjeneste og ere derved førte til stadigt forgæves Forsøg paa at forstaa Forandringerne ud fra det Uforanderlige, det Anskuelige ud fra de rene Begreber. Filosofiens rette Opgave er at forsøge at genvinde den umiddelbare Enhed og den inderlige Sammenhæng i vort Sjæleliv, som det praktiske Liv, i hvis Tjeneste Sproget og Naturvidenskaben staa, atter og atter fører bort fra. Tankelivet maa blive aldeles uinteresseret for at kunne løse denne Opgave, --- for at faa Øje paa, hvad der er det virkeligt Givne, --- Grundlaget for, hvad Sproget søger at udtrykke, og hvad Videnskaben søger at formulere.

Titelen paa det Skrift, med hvilket Bergson først slog igennem, og som i mine Øjne stadigt er hans bedste Skrift --- Les Données immédiates de la Conscience (Paris 1889) --- er i denne Henseende karakteristisk.

2. Henri Bergson er født i Paris i Aaret 1859, og hans vigtigste Skrift udkom altsaa, da han var 30 Aar gammel. Om hans Livsgang foreligger der ikke nærmere Oplysninger; den synes at gaa i Et med Udviklingen af hans Tankeliv. Efterat have gennemgaaet École Normale virkede han som Lærer ved forskellige højere Almenskoler i Provinserne og i Paris, derefter en Tid (1898--1900) ved École Normale og virker nu, fra Aar 1900 af, som Professor ved Collège de France.

Han er en glimrende, ofte blændende Forfatter og skal være en udmærket Taler. Hans Forelæsninger ere overfyldte og synes at være blevne en Modesag for det intelligente Bourgeoisi. Det er dog ikke Noget, han selv har jaget efter. Han gør personlig et stilfærdigt, fint og tiltalende Indtryk, og
% --- p. 4 ---
\opage{4}jeg for mit Vedkommende mindes med Glæde den interessante filosofiske Samtale, jeg havde med ham i Paris for nogle Aar siden. Jeg havde allerede dengang korteligt antydet min Forskel fra ham i en historisk Oversigt, og senere har jeg i min Bog „Den menneskelige Tanke“ drøftet forskellige vigtige Punkter i hans Opfattelse \nm{2}.

Bergson betegner, som det allerede vil ses af det Foregaaende, en Reaktion, baade mod Rationalisme og Realisme, en Reaktion, der idethele kan spores i Nutiden paa forskellige Omraader. Der gives maaske dem, for hvem Dommen over Bergson allerede er fældet med Ordet Reaktion. Men en Reaktion kan være fuldt berettiget, naar nye Synspunkter tiltrænges, eller naar gamle Synspunkter, der en Tid vare skudte tilside, fordi Forskningen paa sine bugtede Veje for en Tid ikke havde Brug for dem, igen godtgøre deres Nødvendighed. Tænkningens Historie viser, at den skrider frem gennem en stadig Brydning mellem Aktioner og Reaktioner, der gensidigt vække og indirekte supplere hverandre. Syntese og Analyse, Kontinuitet og Diskontinuitet, Enhed og Mangfoldighed, Anskuen og Eftertanke, uvilkaarlig Griben og bevidst Søgen --- alle disse forskellige Retninger eller Tendenser i Tankelivet maa stadigt vexelvirke, men dette kan i Virkeligheden kun ske ved, at snart den ene, snart den anden faar Overvægten. Det er snart indre, snart ydre Aarsager, der betinge, hvilken Tendens der for Tiden raader. Snart er der en aandelig Trang i den ene eller anden Retning, der kræver Tilfredsstillelse, snart er det nye Iagttagelser, der kræve andre Former og stille andre Opgaver, end der hidtil gjorde sig gældende. Min Opgave skal i det Følgende være at undersøge, hvorledes det i dennne % PRINTED AS IS: dennne (three n)
Henseende forholder sig med Bergson’s Indlæg i den filosofiske Diskussion.

Han fører sin Sag med en Glans og en Veltalenhed, der have bidraget meget til at vække Opmærksomhed og Tilslutning. Og tillige er der en Strømning i Tiden, som kommer ham imøde; fra saa forskellige Sider som Katolicismen og Syndikalismen er der vist ham Sympati.
% --- p. 5 ---
\opage{5}Det kan strax siges, at Tankernes Begrundelse og nøjagtige Gennemførelse ikke staar paa Højde med Fremstillingens Glans. Og der melder sig strax en Hovedvanskelighed for en Opfattelse som Bergson’s. Naar baade Sproget og det videnskabelige Tankearbejde lide af en radikal Ufuldkommenhed, der udspringer af, at de helt igennem ere formede ved deres Arbejde i praktiske Formaals Tjeneste, hvilke Udtryk og hvilke Tankeformer kan der da raades over til at fremstille og udvikle den højere Art af Forstaaelse, hvortil Bergson vil føre os?

Alle maa indrømme, at Eftertanken stedse maa bygge, og virkeligt bygger, paa Anskuelsens Grundlag, hvad enten man nu ved Anskuelse tænker paa Sansning, paa Erindring eller paa Fantasi. Og man kunde tænke sig, at vor Anskuelsesevne var saa omfattende og saa bevægelig, at den stadigt kunde virke paa fyldestgørende Maade, uden at den analyserende og abstraherende Eftertanke behøvede at træde til. Jeg har i min Bog „Den menneskelige Tanke“ omtalt denne Mulighed, men strax føjet til, at det nu engang ikke forholder sig saaledes. Der kræves atter og atter bevidst Sammenligning og Ordnen af de Emner og Elementer, der udgøre Anskuelsens Indhold, dersom dette skal fremtræde i fuld Klarhed og Begrundelse. Den blinde, uvilkaarlige og ubevidste Maade, paa hvilken Emner og Elementer forbindes i umiddelbar Anskuelse, viser sig erfaringsmæssigt ikke at være tilstrækkelig. Selv om man maa give Goethe Ret i, at man skal dvæle i Anskuen saa længe som muligt og ikke for tidligt lade den opløsende Eftertanke og Dømmen begynde sit Værk, saa vil dog, naar den aandelige Udvikling ikke staar stille, det Punkt komme, hvor dette Værk maa gøres \nm{3}.

Bergson lægger ligesom Goethe meget stor Vægt paa Anskuelsen og betragter Eftertanken som et Tegn paa Ufuldkommenhed. „Dersom vore Sanser og vor Bevidsthed havde en ubegrænset Rækkevidde, --- dersom vor Evne til ydre og indre Iagttagelse var ubegrænset, --- vilde vi aldrig tage vor Tilflugt til Evnen til at danne Begreber og drage Slutninger. At danne Begreber er en nødtvungen Udvej i de Tilfælde, hvor man ikke kan iagttage, og vi nødes kun til at hjælpe os med Slut%
% --- p. 6 ---
\opage{6}ninger, hvor det gælder at udfylde den ydre eller indre Iagttagelses Huller og at udvide deres Omfang. Jeg nægter ikke abstrakte og almene Begrebers Nytte --- ligesaalidt som jeg bestrider Værdien af Bankanvisninger. Men ligesom en Anvisning kun er et Løfte om Guld, saaledes har et Begreb kun Værdi ved de mulige Iagttagelser, det repræsenterer. Derom er, paastaar jeg, Alle enige. Ti Alle indrømme, at de sindrigst byggede Begreber og de lærdest konstruerede Slutningsrækker falde sammen som Korthuse, naar en eneste virkeligt iagttagen Kendsgerning strider mod disse Begreber og Slutninger. Der er desuden ingen Metafysiker og ingen Teolog, som ikke vil være villig til at paastaa, at et fuldkomment Væsen maa kende Alt gennem umiddelbar Anskuen (intuitivement), uden at behøve at gaa gennem Slutningens, Abstraktionens og Generalisationens Mellemled“ \nm{4}.

3. Bergson indrømmer naturligvis, at vi nu engang ikke ere saa fuldkomne Væsener, at vi kunne hjælpe os med blot Anskuen. Og selv om vi i denne Henseende vare fuldkomnere, end vi ere, saa vilde vi dog gerne meddele os til hverandre om, hvad vi skue, og det kan kun ske ved Sproget, der ifølge Bergson er blevet ubrugeligt, ligesom Videnskaben er det. Hverken Begreber eller Billeder slaa til, naar det gælder at gengive de umiddelbare Anskuelsers Indhold; Begreberne staa i et udvortes Forhold til hverandre og udtrykke hvert for sig kun en enkelt Side eller Egenskab ved det Anskuede, og intet Billede er fuldt udtømmende. Billedet har nu alligevel det Fortrin for Begrebet, at det giver noget Individuelt og Konkret, og Bergson opfordrer os derfor til at gøre et Forsøg med Billeder. Forsøg, siger han, at tænke paa dig selv, paa dit eget indre Liv, saaledes som det udfolder sig i Tiden, dine egne skiftende Tilstande! Forskellige Billeder kunne da skyde sig frem. Der er vel intet af dem, der slaar til, men de kunne alle pege i samme Retning. Du finder maaske en Modsætning mellem den ydre Skorpe og de indre Lag. Men dette Billede lader ikke den stadige bevægelige Sammenhæng (la continuité d’écroulement) komme til sin Ret. Eller det er dig, som om Noget viklede sig op inde i dig, ligesom et Nøgle Garn
% --- p. 7 ---
\opage{7}rulles op, eller omvendt, som om Noget rullede sammen i dig. Heller ikke dette Billede slaar til, fordi det, der rulles op eller ind, forestilles som ensartet og som udstrakt i Rummet. Prøv da Billedet af et Spektrum med tusinde Nuancer! Men ogsaa disse Nuancer ligge udenfor hverandre, og desuden ligge de fast; der mangler den Henbølgen, som indre Iagttagelse direkte godtgør. --- Intet Billede slaar altsaa til. Men alle de fremdragne Billeder pege dog i samme Retning. Og netop derved, at hvert Billede halter, men alle Billeder pege i samme Retning, kan der bibringes et Stød, vækkes en Trang til selvvirksom Frembringen af en indre Tilstand eller Holdning (attitude) af den Beskaffenhed, som den Talende vil have den, han taler til, ind i. Kun indirekte kan en aldeles umiddelbar Oplevelse fremkaldes. Det er kun ved Selvvirksomhed, ved en Villiesakt, at vi komme til at staa Ansigt til Ansigt med det mest Umiddelbare og Uvilkaarlige i vort indre Liv, det, som Filosofen vil føre os tilbage til.

Denne Fremstiling, % PRINTED AS IS: Fremstiling (Fremstilling)
der findes i „Indledning til Metafysiken“ \nm{5}, er maaske det Interessanteste, som Bergson har ydet. Han har set, at enhver dybere gaaende aandelig Paavirkning kun kan ske indirekte, og at der derfor hører Kunst til at øve den, ligesom der kræves baade Modtagelighed og Evne til Selvvirksomhed for at kunne vækkes til at finde sig selv, opdage det Inderste i sit indre Liv. Det har Filosofien vidst ligefra Sokrates’ Tid af, men den har ikke altid gjort Brug deraf. Baade Systematiseren og Skoledannelse have her været Hindringer.

Man maa vel bemærke de Betingelser, Bergson hævder for Opnaaelse af den umiddelbare Skuen, der for ham er det Højeste, men som Livets Praxis saa let skjuler for os. Det er ingenlunde, som man undertiden har sagt, Bergson’s Mening at begunstige intellektuel Ladhed. Der skal være unge franske Studerende, som ved Bergson’s Lære have faaet Foragt for videnskabeligt Arbejde. Dette Arbejde kunde, mener de, være nødvendigt for Praktikere, for Ingeniører, Læger o. a., men for Filosofer, hvis Omraade begynder, hvor Videnskaben ender, gælder det om at fordybe sig i de umiddelbare Oplevelser, i
% --- p. 8 ---
\opage{8}den rene Skuen! --- Med Rette har Bergson protesteret mod denne Tydning af hans Filosofi \nm{6}. Selv hvor han tænker paa det „umiddelbart Givne“, som maa være tilgængeligt for Alle, naar de blot arbejde sig ud af det praktiske Livs Synsmaader og Sprogets udstykkende Mekanik, selv dér hører der, som allerede bemærket, baade Modtagelighed og Anstrengelse til for at kunne vende tilbage til de umiddelbare Oplevelser, ligesom der jo ogsaa er særegne Vanskeligheder at overvinde, naar man vil udtrykke, hvad man saaledes oplever. Man kan altsaa ikke drømme sig ind i „det umiddelbart Givne.“

Nu skelner Bergson, som det i det Følgende vil vise sig, ikke tydeligt mellem „det umiddelbart Givne“ i den omtalte Betydning og den Intuition, den umiddelbare Skuen, som er Et med en Helhedsopfattelse, der kan staa som Afslutning eller Afrunding af videnskabeligt Tankearbejde. Forsaavidt han taler om en Intuition af denne sidste Art, er det selvfølgeligt, at den kun kan naas gennem Eftertankens Skærsild. Men der er naturligvis altid dem, der vilde forsøge paa at komme ind i Paradiset uden at gaa igennem Skærsilden.

Det er, efter min Opfattelse, som den følgende Fremstilling skal søge at begrunde, snarere en Art kunstnerisk Skuen, Bergson baner Vej for, end en højere Videnskab. En saadan Videnskab, den være nok saa høj, behøver dog stadigt Begreber og Slutninger, Billeder og Sammenligninger. Hvem borger for Gyldigheden af den Intuition, der skal kunne opstaa, efter at man har gennemgaaet Forskningens Skærsild? At der ad den indirekte Vej, Bergson saa godt beskriver, kan vækkes en indre Bevægelse, en Fordybelse og en Levendegørelse indenfor Sindets Verden, er psykologisk muligt; men hvori ligger det, at der derved skal kunne fremtræde Noget, der, hvad Gyldighed, hvad Sammenhæng med Tilværelsens Inderste angaar, staar over Videnskabens Synsmaader og Resultater? ---

4. Ved sin stærke Fremhæven af det umiddelbart Givne og den umiddelbare Leven og Skuen, i Modsætning til Eftertanken, dens Former og dens Resultater, minder Henri Bergson om en Række af tidligere franske Forfattere, uden at han
% --- p. 9 ---
\opage{9}kan siges at være paavirket af dem. Han tilhører en Type, der atter og atter har gjort sig gældende i fransk Aandsliv.

Jeg nævner først Montaigne. Bag hans Skepsis staar en Hengivelse til Naturen, „vor store Moder“, der umiddelbart rører sig i os alle, i Enhver efter hans Ejendommelighed. Og Naturen i os og over os kan ikke udtrykkes ved Hjælp af bestemt afgrænsede Begreber. --- Pascal stiller „Hjertet“s umiddelbare Erfaring i Modsætning til abstrakt Bevisførelse, „esprit de finesse“ i Modsætning til „esprit de géométrie“: hin ser Tingen i dens Ejendommelighed paa én Gang og med et eneste Blik, denne fatter den gennem en Række af Slutninger. --- En Analogi hertil frembyder Rousseau’s Modsætning mellem Natur og Kultur, Følelse og Forstand, Samling og Spredthed. --- Maine de Biran fremhæver forskellige Trin af umiddelbar Leven: vie animale, vie humaine og vie de l’esprit, og han dvæler særligt ved „det indre Menneske“, som ikke kan ytre sig udad saaledes, at fuld Forstaaelse bliver mulig. Fra Maine de Biran gaar der en Udviklingslinie gennem Félix Ravaisson, Lachelier og Boutroux til Bergson. --- Ogsaa med Guyau frembyder Bergson Analogi. Guyau’s Hovedtanke er Udstraaling og Udformning (expansion) af Sindets Kræfter. Tanken spiller i Forhold hertil en underordnet Rolle. Og i sit Digt „Le mal du poète“ skildrer han Digterens Sygdom som bestaaende i, at intet Billede formaar udtømmende at udtrykke Oplevelsens Fylde.

En tilsvarende Betydning havde i den tyske Filosofis Historie Jacobi, Herder og Hamann overfor Kant ved Overgangen til det nittende Aarhundrede. Ogsaa de hævdede Aandens samlede, koncentrerede Virken i Modsætning til Analyse og Kritik, Tingenes levende, konkrete Sammenhæng i Modsætning til de skarpe Distinktioner i den kritiske Filosofi. Og i en senere Tid taler Friedrich Nietzsche i overstrømmende Lidenskab den mægtige Livsudfolders Sag overfor Rationalismen \nm{7}. ---

Saadanne Aandsretninger have for Filosofien den Betydning, at de gøre opmærksom paa, hvilket stort Indhold der stedse vil blive tilbage, efterat Tanken har gjort sit Bedste med at analysere, skelne og sammenfatte de Emner, den sysler
% --- p. 10 ---
\opage{10}med. Tanken øser stedse ud af en Fylde, den ikke vil kunne udtømme. Dette vil dog ikke kue dens Mod eller forringe dens Betydning; tvertimod godtgør det, at den aldrig vil mangle Opgaver.

De største Skikkelser i fransk Filosofi, Descartes og Comte, høre ikke til denne Type. Filosofiens sidste Ord maa være Tankens Ret og Myndighed, som ikke behøver at medføre Miskendelse af Livets Mangfoldighed. I Nutiden repræsenteres Comte’s positivistiske Retning med Energi af Sociologen Durkheim. Som Repræsentanter for Descartes’ rationelle Retning kunne nævnes, foruden Émile Boutroux, to nyligt bortgangne Forskere, Alfred Fouillée og Matematikeren og Filosofen Henri Poincaré. Henri Bergson staar i bestemt Modsætning til begge Retninger. Vi skulle nu i det Følgende nærmere undersøge, ud fra hvilke Synspunkter han imødegaar dem.

\chapterend

% --- p. 11 ---
\chapterhead{INTUITION}

\opage{11}5. \initial{D}{er} er ifølge Bergson en stor Modsætning mellem, hvad der umiddelbart og uvilkaarligt fremtræder i og for Bevidstheden, og hvad Eftertanken skelner, sondrer og saa bagefter søger at forbinde. Forud for og ud over al Sondren og Forbinden gaar en stadig Strøm af Oplevelser, af indre Forandringer i uafbrudt, levende Sammenhæng, som danner det nødvendige Grundlag, paa hvilket Eftertanken kan udfolde sig, og tillige staar som det Ideal, til hvilket vi fra Eftertankens Arbejde søge tilbage, naar Tankearbejdet skal afsluttes.

Intuitionen, hvorved Bergson mener Fordybelse i denne umiddelbare, strømmende Sammenhæng, er for ham baade Grundlaget for og Afslutningen af Tankearbejdet. Og jeg kan ikke se Andet, end at der er en afgørende Uklarhed hos ham med Hensyn til Forholdet mellem Intuition som psykologisk Forudsætning og Intuition som tankemæssig Afslutning, en Uklarhed, der tildels har voldet de Mistydninger, hans Filosofi har været udsat for.

Ligefra sin første Optræden stiller Bergson sig i berettiget Modsætning til en Retning, som i den forudgaaende Menneskealder øvede en stor Indflydelse paa det filosofiske, særligt det psykologiske Omraade. Det er den Retning, man plejer at betegne som den engelske Associationspsykologi, og som betegnes ved Navnene Hume, James Mill, Stuart Mill og Herbert Spencer. Denne Retning har havt den store Fortjeneste at have banet Vejen for en mere videnskabelig Psykologi ved at opfatte Sjælelivet som bestaaende af en Mængde Elementer --- Fornemmelser,
% --- p. 12 ---
\opage{12}Forestillinger, Følelser og Drifter --- der vare selvstændige i Forhold til hverandre og kun paa udvortes Maade, nemlig derved, at de ofte forekom sammen, kunde komme i fastere Forbindelse indbyrdes. Sjælelivet var et Kaos, der kun ved Vanens Magt kom til at danne et ordnet Hele. Denne Teori har bidraget meget til at føre ud over vage Beskrivelser og mystiske Henvisninger til Sjælelivets Enhed. Den har gjort det lettere at stille bestemte psykologiske Spørgsmaal. Og den har bidraget meget til at bringe Psykologien i nøjere Forbindelse med Naturvidenskaben. Men den forvexlede afgjort sine teoretiske og abstrakte Synspunkter med, hvad Iagttagelse af det virkeligt givne Sjæleliv udviser, og den kom særligt i Strid med den Kendsgerning, at intet sjæleligt Element kan forstaas i sin Ejendommelighed uden gennem den Sammenhæng og Vexelvirkning, i hvilken det staar med andre sjælelige Elementer. Elementerne som saadanne ere overhovedet stedse Produkter af Analyse, af Eftertanke og maa ikke forvexles med, hvad det umiddelbart givne Sjæleliv frembyder.

Baade i og udenfor England blev denne Opfattelse imødegaaet. Her er saaledes at nævne Wundt’s „Physiologische Psychologie“ (1874) med dens stærke Betoning af Opmærksomhedens („Apperceptionens“) Betydning for Sjælelivets Udvikling. Samme Aar udkom Thomas Hill Green’s „Introduction to Hume“ med en skarpsindig Kritik af de Forudsætninger, hvorfra „Associationspsykologiens“ berømteste Repræsentant var gaaet ud. Samme Aar udkom ogsaa min Bog „Den engelske Filosofi i vor Tid“, i hvilken de kritiske Bemærkninger allerede antydede den Opfattelse, der gennemførtes i min „Psykologi paa Grundlag af Erfaring“ (1882), særligt ved Hjælp af Forholdet mellem sundt og sygt Sjæleliv. I England betegnede Francis Herbert Bradley’s „Ethical Studies“ (1876) det afgørende Brud med den traditionelle engelske Skole \nm{8}. Paa den beskrivende Psykologis Omraade gjorde William James’ Afhandling om Tankestrømmen (The Stream of Thought. Mind 1884, senere optrykt som Kap. 9 i hans „Principles of Psychology“) Epoke ved en glimrende Skildring af den store, uvilkaarlige Sammenhæng i Sjælelivet, der er Grundlaget for al Selviagttagelse og al psykologisk Analyse.
% --- p. 13 ---
\opage{13}Bergson er altsaa ikke den Første, der har kritiseret Overførelsen af rent mekaniske Analogier paa Sjælelivets Omraade og søgt at lade det virkeligt Givne komme mere til dets Ret. Men han har gjort det med en Kraft og med en Fremstillingskunst, der langt overgaar hans Forgængeres, i sit Hovedværk „Les Données immédiates de la Conscience“ (1889), hvis Tankegang nu skal gengives.

6. Vore sjælelige Oplevelser foregaa --- hvad vi ville opdage, naar vi frigøre os for de rumlige og mekaniske Forestillinger, der fra Naturvidenskaben have trængt sig ind i Aandsvidenskaben --- i Lighed med henglidende, sammenhængende Strømme. De foregaa i Tiden, ikke i Rummet. Den ene Tilstand gaar umiddelbart over i den anden, maaske netop naar denne er af helt modsat Art. Det er ingen ensartet Strøm; der gør sig stadigt kvalitative Forandringer gældende. Farvefornemmelse afløser Farvefornemmelse, Sorg afløser Glæde, Opmærksomhed Sløvhed, Spænding Hvile, det Store det Smaa o. s. v. Og disse Oplevelser forstaas hver for sig kun gennem deres Sammenhæng. De ere ikke skilte fra hverandre paa samme Maade som Tingene i den ydre Verden, der hver befinde sig paa sit Sted i Rummet.

En Følge heraf er, at sjælelige Oplevelser ikke umiddelbart kunne maales, saaledes som det, der foregaar i Rummet. Vi kunne ikke afgøre, om to Fornemmelser ere absolut lige stærke. Vi kunne ikke addere og subtrahere Fornemmelser. Naar vi sige, at vi maale dem, sker Maalingen gennem deres Forhold til de ydre Paavirkninger; saaledes naar vi maale Varmefornemmelserne ved Hjælp af Termometret. Størrelsesbegrebet hører overhovedet ikke hjemme i Aandsvidenskaben, men i Videnskaben om de materielle Fænomener. Under Kampen for Livets fysiske Opretholdelse er det af Vigtighed at kunne tælle, maale og veje, og det er Naturvidenskabens Storhed, at den formaar at gennemføre en kvantitativ Synsmaade for hele den materielle Verden. Men alt det Arbejde, der her øves, er fra først af fremkaldt af Livets praktiske Behov og staar stadigt i dettes Tjeneste.
% --- p. 14 ---
\opage{14}Det Forsøg, Fechner har gjort paa at maale sjælelige Fænomener, bliver som Følge af denne Betragtningsmaade Genstand for Kritik fra Bergson’s Side. Han holder sig dog herved udelukkende til Vanskeligheden ved at finde psykiske Enere, det vil sige, saa simple Fornemmelser, at de med absolut Nøjagtighed kunne sættes istedetfor hverandre, og at mindre simple Fornemmelser skulde kunne betragtes som Summer. Denne Kritik er ikke ny; den blev --- som Bergson selv nævner --- allerede 1875 fremført af Matematikeren Jules Tannery \nm{9}, og dens Rigtighed er almindeligt anerkendt. Bergson udtrykker træffende dens Hovedtanke i følgende Ord: „Det væsentligste, efter min Mening det eneste Spørgsmaal, er, om en Kontrast mellem A og B virkeligt kan være ligestor med en Kontrast mellem B og C, naar de to Kontraster finde Sted mellem forskellige Elementer. Saasnart man har kunnet fastslaa, at to Fornemmelser kunne være ligestore uden at være identiske, vil Psykofysiken være grundlagt.“

Der er andre Sider ved Fechners psykologiske Undersøgelser, som Bergson lader ligge, uagtet de netop vilde passe godt ind i hans Teori. Det er den Maade, paa hvilken Fechner fører et experimentalt Bevis for, at en Fornemmelse ikke har nogen absolut Størrelse, der er den samme under alle Forhold, men at den Størrelse eller Styrke, den optræder med, beror paa dens (eller rettere den tilsvarende Paavirknings) Forhold til forudgaaende eller samtidige Paavirkninger.

Desuden kunde der være henvist til, at alle Maalinger ogsaa paa det materielle Omraade bero paa Analogi, Overførelse af Talforhold paa Forholdet mellem Rumsfænomener, og at det heller ikke her er muligt at paavise Enere, der ere absolut uforanderlige under alle Forhold.

7. En betydningsfuld Ejendommelighed ved alle sjælelige Fænomener, naar de fremtræde i deres fulde Umiddelbarhed, er, at de udfolde sig i Tiden. Tidsbegrebet spiller da ogsaa en Hovedrolle i Bergson’s Filosofi. Nu er dette Begreb jo ogsaa et vigtigt Begreb i Naturvidenskaben, der mere og mere bliver til en Lære om Bevægelser. Men her finder Bergson netop
% --- p. 15 ---
\opage{15}det mest skæbnesvangre Punkt, hvor uberettigede materielle Analogier fra Naturvidenskaben ere trængte ind i Aandsvidenskaben. Naturvidenskabeligt opløser man --- mener han --- Tiden i Øjeblikke, og man opfatter dernæst hvert Øjeblik som et Sted eller en lille Strækning. To Tider kaldes i Naturvidenskaben ligestore, naar ens Legemer under ens Forhold have gennemløbet ligestore Strækninger. Man stiller da Rummet istedetfor Tiden. Eller rettere, den psykologiske Tid afløses af den mekaniske eller geometriske. I den psykologiske Tid, den Tid, i hvilken vi virkeligt leve, og i hvilken vi opleve, hvad der er vort Livs virkelige Indhold, har ethvert Øjeblik sin ejendommelige Beskaffenhed, er udfyldt paa sin særegne Maade, og intet Øjeblik er derfor ganske ligt det andet. --- Den psykologiske Tid kalder Bergson la durée, nøjagtigere la durée qualité, til Forskel fra Naturvidenskabens temps matérialisé, der kan symboliseres ved en Strækning i Rummet.

Tendensen til at bruge Rumsbilleder, ogsaa naar Talen er om vore indre Oplevelser, stammer ifølge Bergson ikke blot fra det praktiske Arbejdes og Naturvidenskabens Indflydelse, men ogsaa fra Trangen til Meddelelse. Rummet yder os de bedste Tegn og Symboler, naar vi ville gøre os forstaaelige for hverandre. „Vi have en uvilkaarlig Tendens til at gøre vore Oplevelser faste og solide for at kunne udtrykke dem ved Sproget. Deraf kommer det, at vi sammenblande selve vor Følelse, der er i en stadig Vorden, med dens ydre Genstand, der har en stadig Væren, og ganske særligt med Ordet, som betegner denne Genstand.“ I Rummet have vi en Form, der baade er fælles for os alle og er maalelig og klart bestemmelig. I Anskuelsen af et ensartet Rum er man allerede paa Vej til det sociale Liv. \nm{10} Saa har man nemlig et fælles Skema, ved hvis Hjælp man kan ordne og klare sine Meddelelser. Alt kan finde sin Plads der, naar man ikke tager det saa nøje med Forskellen mellem det umiddelbart Givne og dets symbolske Udtryk. Hvis vi derimod netop ville forsøge at faa Øje paa, hvad der virkeligt er givet, saa maa vi stræbe at se bort fra alle kvantitative Bestemmelser og fra alle mulige Symboler. ---
% --- p. 16 ---
\opage{16}8. Der er i hele denne Tankegang, som Bergson har fremstillet paa en glimrende Maade, meget Berettiget, og Bergson har stor Fortjeneste af ved sin indtrængende Skildring af sjælelige Tilstandes Ejendommeligheder at have stillet deres Forskel fra andre Sider i vor Erfaring i et klarere Lys. Nogle kritiske Bemærkninger maa dog gøres, inden vi gaa videre.

Til de kvantitative Begreber, der maa holdes borte fra Anvendelse paa Sjælelivet, regner Bergson ogsaa Talbegrebet. Det hviler nemlig efter hans Opfattelse paa Diskontinuitet mellem Steder i Rummet. Han mener, at vi kun kunne tælle, hvad der kan henføres til forskellige Steder. Men man kan jo meget godt tælle Kvalitetsforskelligheder. Jeg kan tælle mine Tanker og mine skiftende Stemninger indenfor en vis kort Tid. Holbergs „Vægelsindede“ havde femten Sind udi én Time. Talbegrebet forudsætter ganske vist Diskontinuitet, men ikke nødvendigvis Diskontinuitet mellem Steder i Rummet. Efter Bergson’s egen Opfattelse er jo det umiddelbart Givne ingenlunde ensartet og uforanderligt; det frembyder stedse nye Forskelligheder, og her er da et tilstrækkeligt Grundlag for Talbegrebet. Man kunde meget vel tænke sig Væsener, der ingen Rumsforhold kendte, og som dog kunde have dannet Talbegrebet. Tilsidst er det i Grunden altid vore egne Tanker eller Opmærksomhedsakter, vi tælle. \nm{11}

Det her berørte Punkt hænger sammen med, at Bergson overhovedet gør for stor og udvortes Forskel mellem det umiddelbart Givne (hvad jeg plejer at kalde Emnerne) og det psykiske Arbejde. Intet er os givet, intet Emne fremtræder for os uden psykisk Arbejde, hvad enten vi saa mærke dette Arbejde eller ikke. Naar man stiller et umiddelbart Givet (et Emne) i Modsætning til den Virksomhed, ved hvilken det kan blive ændret gennem Analyse og Abstraktion, saa udelukkes derved ikke, at selve hint Givne for en Del kan skyldes en forudgaaende psykisk Virksomhed. Det forudsætter Opmærksomhed for overhovedet at blive til. Det har samlet sig til en vis Helhed for os i Kraft af den samme Proces, der med større Grad af Bevidsthed gaar for sig, naar vi sammenarbejde spredte Iagttagelser til et Hele. Den forskellige Styrke, hvormed enkelte Elementer i det Givne
% --- p. 17 ---
\opage{17}fremtræde, skyldes dels Kontrastforhold, dels forskellig Følelsesbetoning. Under Indflydelse af saadanne Kræfter sker det, at vore Sansebilleder, vore Erindringer og vore Fantasier uvilkaarligt skifte Karakter, saa at vi maaske først længe efter blive os disse Forandringer bevidst. Det er derfor, jeg i min Bog om „den menneskelige Tanke“ har sagt, at hvad der i hvert enkelt Tilfælde kaldes umiddelbart givet, kun fortjener dette Navn i Forhold til et vist bestemt Arbejde, medens det fra andre Synspunkter og i Sammenligning med andre Emner maaske vilde være at betegne som et Arbejde.

Hvis der var den Modsætning mellem Emne og Eftertanke, som Bergson hævder, vilde det være umuligt, at Eftertanke overhovedet vilde kunne opstaa, og det vilde desuden blive umuligt at give nogen Beskrivelse af det umiddelbart Givne. Det vilde være, som om der i os boede to Personer, af hvilke den ene havde Emnerne, den anden Eftertanken; de vilde ikke forstaa, eller endog blot opdage hinanden. ---

9. Det er for Bergson, som om et Syndefald havde fundet Sted. Forledet af de Slanger, Livets Praxis udklækker, have vi vendt det umiddelbart Givne Ryggen og hengivet os til den udstykkende og abstraherende Eftertanke.

Spørge vi nu, naar dette Syndefald er foregaaet, henviser Bergson os til den Bog, han selv betragter som sit Hovedværk: \textit{Évolution Créatrice} (1907), og vi føres derved fra Psykologien over til Biologien.

Syndefaldet skete, da Forstanden (l’intelligence) afløste Instinktet. Instinktet staar Livet nærmere, end Forstanden gør. Det er en umiddelbar Fortsættelse af Væxten, af Organisationsprocessen. Naar den lille Kylling knuser sin Æggeskal ved et Stød med Næbbet, handler den af Instinkt, og dog gør den ikke Andet end at fortsætte den Bevægelse, der har ført den gennem Fostertilstanden. Instinktet udspringer af selve Livsvirksomheden. Forstanden opstaar kun, hvor Instinktet viser sig utilstrækkeligt, --- som et Middel til at danne Redskaber. Den er ligesaa udvortes og mekanisk, som Instinktet er organisk. Forstanden bevæger sig uden om sin Genstand, søger at finde
% --- p. 18 ---
\opage{18}saa mange forskellige ydre Synspunkter for den som muligt, men formaar ikke at gøre sig til Et med den. Den drager Genstanden til sig, istedetfor at trænge ind i den. Instinkt og Forstand staa i den skarpeste Modsætning til hinanden. En Mængde afledede Instinkter vil Forstanden maaske kunne give en videnskabelig Forklaring af, men en fuldstændig Analyse af Instinktets Væsen vil den aldrig kunne give; dertil er den for formel. Instinktet kan ikke udtrykkes i Forstandsbegreber. Det er Kontinuitet og Diskontinuitet, der her staa overfor hinanden. Der er i Verden Mangt og Meget, som Forstanden ene er istand til at søge, men som den ikke vilde kunne finde. Instinktet vilde kunne finde det, men vilde aldrig søge derefter!

Dersom den Energi, der fører til Instinkthandlingen, kunde vaagne til Bevidsthed, inderliggøres til Erkendelse istedetfor at omsættes til Praxis, --- dersom den potentielle eller slumrende Bevidsthed i Instinktet kunde vaagne til fuld Klarhed, --- saa vilde Livets inderste Hemmeligheder blive aabenbarede. Men Livets Krav medførte Forstandens Udvikling, for at Midler og Redskaber kunde skaffes tilveje, og dermed var det skæbnesvangre Skridt gjort. \nm{12} ---

Denne Bergson’s Lære om Forholdet mellem Instinkt og Forstand fremkalder uvilkaarligt den Bemærkning, at der dog ikke er nogen Grund til at vente sig større Oplysninger om Livets Gaade af Instinktet end af Forstanden. Hvis Bergson har Ret i, at Forstanden er praktisk, saa gælder det dog i nok saa høj Grad om Instinktet. Ti praktisk er det jo dog at gaa lige til Maalet, hvor det kan lade sig gøre, istedetfor at gaa unødvendige Omveje. Instinktet og Forstanden ere to Midler til at naa et Maal, og de svare hver for sig til sine bestemte Livsforhold. Der er ikke Grund til at søge noget særligt Mystisk i Instinktet. Det maa vel antages, at der i Instinktet rører sig en Trang, men netop fordi „Bevidstheden sover“, er det en blind Trang, der ingen Hemmeligheder kan aabenbare os.

Bergson antager en vis Analogi mellem Instinkt og Intuition. Intuitionen er „et uinteresseret Instinkt“. \nm{13} Men fordi Intuitionen kan kaldes et Instinkt, idet den virker uvilkaarligt, kan omvendt ikke Instinktet betragtes som en Intuition, og det er kun
% --- p. 19 ---
\opage{19}derved, at det faar sit mystiske Præg. I sig selv træller det nok saa meget for Livets Praxis som „Forstanden“. Og naar Instinktet skal kunne blive uinteresseret, emancipere sig fra Trældom i praktiske Formaals Tjeneste, hvorfor skulde saa ikke ogsaa „Forstanden“ kunne det? Er desuden ikke ogsaa Tankearbejdet en Art af Liv, der, hvis det kunde vaagne til Bevidsthed om sig selv og sine Love, ogsaa kunde kaste Lys over Livets Væsen og Vilkaar? Livet rører sig paa alle Trin --- i Tanken saavel som i Instinktet og i Organisationsprocessen. Og overalt gaar den uvilkaarlige Virken og Stræben forud for Besindelse og klar Bevidsthed.

Ligesaalidt som det var begrundet og berettiget at stille det i psykologisk Henseende „umiddelbart Givne“ i absolut Modsætning til Eftertanken, ligesaalidt er det altsaa muligt at stille Instinktet i absolut Modsætning til Forstanden og betragte dem som staaende i højst forskellig Afstand fra Livets inderste Væsen.

10. Spørges der, om der ikke gives nogen Oprejsning efter Syndefaldet, svarer Bergson, at det gælder om at arbejde sig ud af den stadige Optagethed af praktiske Interesser, der medfører Opgaaen i det Ydre og Materielle. I sin „Indledning til Metafysiken“ havde han allerede stillet Opfordring i denne Retning, og i sit naturfilosofiske Værk om „skabende Udvikling“ gentog han denne Opfordring med særligt Hensyn til Forholdet mellem Instinkt og Forstand. Det koster Villiesanstrengelse at følge denne Opfordring. Ti Livets praktiske Behov, i hvis Tjeneste Forstanden staar, har i den Grad bestemt vor Maade at tænke paa, vor Forstands hele naturlige Tendens (la pente naturelle de l’intelligence), at vi arbejde mod Strømmen, naar vi søge at komme ud over den udvortes, mekaniske Maade at betragte Tilværelsen paa. Vi maa ved en Villiesakt bryde med de videnskabelige Vaner, med selve Tankens Grundfordringer, hvis vi ville naa den Udvidelse af vor Forstaaelse, den Skuen, i hvilken vi kunne komme Ansigt til Ansigt med Tilværelsens inderste Væsen, den Skuen, som det er Filosofiens egentlige Opgave at føre os til.

Dog staar „Forstanden“ og dens Maade at tænke paa ikke
% --- p. 20 ---
\opage{20}blot som en Hindring for Opstigningen til den Intuition, der er den højeste Opgave. Hvis Forstanden ikke var vaagnet og havde gjort sit Arbejde, saa var Udviklingen ikke kommen ud over det blotte Instinkt, der udelukkende og umiddelbart fører til Bevægelse i en bestemt Retning. Forstandens Bevægelighed og den Orienteren i Omverdenen, den fører til, udvikler Kræfter, der nu kunne vendes indad og anvendes til at genopvække de Muligheder for umiddelbar Skuen, der i Instinktet ikke kom ud over den slumrende Tilstand, fordi Bevægelsestrangen strax rørte sig. Nu kan Instinktet blive frigjort for denne Trang, blive helt uinteresseret, og det bliver muligt med intellektuel Sympati at fordybe sig i og føle sig Et med den Livsstrøm, som bevæger sig i os og gennem den hele Tilværelse, af hvilken vi ere Led. Det er en større Verden, vi maa føres ind i, en Verden, Forstanden paa Grund af sine begrænsede Opgaver maatte isolere sig fra. \nm{14}

Om Karakteren og Gyldigheden af den Metafysik, Bergson her ser Muligheden af, vil jeg senere komme til at tale. Her gjaldt det især om at vise, at Adgangen til den Intuition, i hvilken Metafysiken ifølge ham egentligt bestaar, ikke er let, og at den kræver Forstandsarbejdets Skærsild som sin Forudsætning.

Ogsaa fra et erkendelsesteoretisk Synspunkt har han indskærpet, hvad der her biologisk set viste sig som en Opstigen fra Instinkt gennem Forstand til Intuition. I en Forhandling i det franske filosofiske Selskab om Betydningen af Ordet „Intuition“ har han gjort følgende Opfattelse gældende:

„Intuitionen er ganske vist en oprindelig Aandsvirksomhed, der ikke kan udledes af den brudstykkeagtige og udvortes Erkendelse, ved hvilken vor Forstand, som den sædvanligt bruges, udefra anlægger en Række af Synspunkter paa Tingene. Men man maa ikke miskende, at den intuitive Maade at gribe Virkeligheden paa ikke længere er os naturlig under vor Tænknings nuværende Tilstand. Vi maa for at opnaa den fordetmeste forberede os gennem en lang og omhyggelig Analyse og gøre os fortrolige med Alt, hvad der angaar Genstanden for vort Studium. Denne Forberedelse er især nødvendig, naar det drejer
% --- p. 21 ---
\opage{21}sig om almene og sammensatte Virkeligheder som Livet, Instinktet, Udviklingen: videnskabelig og nøjagtig Kundskab om Kendsgerningerne er den foreløbige Betingelse for den metafysiske Intuition, der skal trænge ind i deres Princip.“ \nm{15}

11. Det er klart, at der er en stor Forskel mellem den Intuition, man naar efter at være gaaet gennem Erfaringens og Analysens Arbejde, og den „oprindelige eller umiddelbare Intuition“, som er det umiddelbart Givne, og som frembyder en udelt Kontinuitet, der siden udstykkes gennem Forstandsarbejdet. \nm{16} Det er et Punkt, der staar utydeligt hos Bergson, om det ikke skulde være muligt at vende tilbage til denne oprindelige Intuition uden at gaa Forstandens lange Vej, --- om ikke et Menneske midt ude i Livet kunde frigøre sig fra det praktiske Livs Interesser og Vaner og genvinde det umiddelbare Blik paa sig selv og Tingene. Kunstneren kan det ifølge Bergson, som jeg senere nærmere skal omtale; han ser Tingene i deres umiddelbare Fylde og Individualitet. Men skulde dette ikke være muligt for Enhver, da der jo dog hos Enhver rører sig Noget af det, der i stor Stil rører sig hos Kunstneren? --- Det vilde have været til stor Gavn for Klarheden, hvis Bergson havde udtalt sig bestemt og klart om Forholdet mellem den Intuition, der danner den nødvendige Forudsætning for alt psykisk Arbejde (eller snarere: er dettes første Resultat), og den, der skal kunne danne Højdepunkt og Afslutning for Tankearbejdet som en højere Enhed af Instinkt og Forstand. Da Bergson bruger Ordet Intuition om begge Arter, har Misforstaaelse ligget nær.

Den metafysiske Intuition skal være en Genfødelse af Instinktet. Men det er da en Genfødelse, der ganske forandrer Instinktets Karakter. Instinktet er, som vi have set, nok saa „praktisk“ som Forstanden, men nu skal det være helt uinteresseret. Og medens det er karakteristisk for Instinktet, at det finder uden at søge (i Modsætning til Forstanden, der søger uden at finde), skal nu dets Finden staa som Frugten af en Søgen, skønt det kun er ved en Villiesakt, ved en voldsom Anstrengelse, et Spring, at man tilsidst naar op til den afsluttende Skuen.
% --- p. 22 ---
\opage{22}Fra den umiddelbare Anskuen af Bevægelse og Forandring, som vi i vore Analyser og Abstraktioner forlade, er denne afsluttende Skuen forskellig ved en Uddybelse og en Udvidelse; med fuld Bevidsthed koncentrerer Filosofen sig i Anskuelsen af den kontinuerlige Proces, hvorpaa han har det bedste og det nærmeste Exempel i, hvad han erfarer i sit eget Indre, hvor Fortid og Nutid umiddelbart række hinanden Haanden.

Det er Bergson’s Bestræbelse at fremdrage Sjælelivets umiddelbare og uvilkaarlige Tilstande. Nu kunne disse indbyrdes være højst forskellige; men ved deres fælles Karakter komme de for Bergson til at staa hverandre nærmere end de i Virkeligheden gøre. \nm{17} Den uvilkaarlige og umiddelbare Anskuen, saaledes som den fremtræder i Sansning, Erindring og Tænkning, --- den uvilkaarlige og blinde Forbindelse af Trang og Evne, man kalder Instinkt, --- og den Skuen, der kan danne Afslutning af et strengt og indviklet Tankearbejde, --- det er, trods den Analogi, der er imellem dem, tre højst forskellige aandelige Fænomener, og det fører til mislige Konsekvenser, naar der ikke sondres tilstrækkeligt imellem dem. Det er en vis romantisk Interesse, den samme, der førte Rousseau til at søge fra Kultur til Natur, som fører Bergson til at opsøge og dvæle ved hine tre, indbyrdes meget forskellige Fænomener.

12. Der fremtræder i Tænkningens Historie forskellige Betydninger eller Former af Intuition, og det vil have sin Interesse at se, om Bergson’s Intuition falder sammen med nogen af disse, eller om den, i modsat Fald, betegner en ny og berettiget Form.

a. I Sansning, i Erindring og i Fantasi raader en umiddelbar Anskuen. Saasnart jeg lukker Øjnene op, faar jeg et Sansebillede, som afrunder sig til en vis Helhed. I Erindringen dukke tidligere Oplevelser op som uvilkaarlige Helhedsbilleder. Og i Fantasien sker der en Nydannelse af saadanne Billeder, hvis Tilblivelsesmaade kan være ligesaa ubevidst som den Maade, paa hvilken et Sanse- eller et Erindringsbillede bliver til. Der vil ikke fra først af gøre sig nogen skarp Forskel gældende mellem Sansning, Erindring og Fantasi, og de ville paa mange
% --- p. 23 ---
\opage{23}Maader sammensmelte og forbindes med hverandre. Det er først den kritiske Eftertanke, der fører til skarp Sondring mellem disse tre Arter af umiddelbar Anskuen, idet Tvivlstilfælde nøde til at søge bestemte Kriterier for at undgaa Forvexling af Iagttagelse, Erindring og Fantasi.

Intuition i denne Betydning angaar noget Individuelt og Specielt og staar i tydelig Modsætning til den sondrende og opmærksomme Virksomhed, vi kalde Analyse, saavel som til Abstraktion og Generalisation og til Uddragen af Slutninger.

b. Intuition kan ogsaa betyde et Skøn, en uvilkaarlig Afgørelse. Her betyder Intuition ikke umiddelbar Anskuelse af et individuelt Hele, som i Sansning, Erindring og Tænkning, men en Forvisning eller Overbevisning, der bliver til ved, hvad man har kaldet „naturlig Indskydelse“ eller „naturlig Magi“. Saaledes, naar visse Principer synes at følge „af sig selv“, eller naar man er umiddelbart forvisset om Virkeligheden af, hvad man ser for sig, eller nærer en uvilkaarlig Tillid til en Autoritet. Intuition i denne Betydning af Ordet, som er meget hyppig i daglig Tale, stilledes af Thomas Reid overfor Hume’s Kritik, og af Jacobi og Hamann mod Spinoza’s og Kant’s Filosofi. Ved Dannelsen af Menneskers Livsanskuelse og Tro spiller denne Art Intuition en væsentlig Rolle. Den er en uvilkaarlig Syntese af Oplevelser og Iagttagelser.

Medens Bergson lægger stor Vægt paa den første Art Intuition (a), kommer han ikke særligt ind paa denne anden Art Intuition (b), uagtet den, ikke mindre end Instinkt og Forstand, gør sig gældende i det praktiske Liv, og uagtet den, ikke mindre end de af Bergson beskrevne Arter af Intuition, tjener til at føre ud over den Splittelse, som Eftertanken kan volde. Og i Grunden nærmer hans egen „Intuition“ sig stærkt til denne Art.

c. En tredie Art af Intuition have vi i den umiddelbare Opfattelse af Forskel eller Identitet mellem to samtidigt eller successivt fremtrædende Sansnings-, Erindrings- eller Fantasibilleder. Det er en Tankeakt, som virker i al Sammenligning, og som fører til, at to Fornemmelser eller Forestillinger enten holdes ude fra hinanden eller betragtes som Et, saaledes at de
% --- p. 24 ---
\opage{24}for Fremtiden kunne sættes istedetfor hinanden. Intuition betyder her umiddelbar Erkendelse af et Forhold. Det er Eftertankens enkelte Skridt, som ved at gentages og fortsættes kan føre til sammenhængende Forstaaelse og Oversigt.

Intuition af denne Art staar aabenbart ikke i Modsætning til Forstand, men betegner netop en Form for Forstandsvirksomhed. Og den staar kun forsaavidt i Modsætning til Intuition af første Art (a), som den angaar Forholdet mellem Billeder, ikke selve Billederne.

Denne Art Intuition er beskreven af Descartes (i hans „Règles pour la direction de l’esprit“) og senere af Henri Poincaré (i „Science et Méthode“). Den betegner hos dem selve Overgangen fra Anskuen til Analyse eller til Bevis, en Overgang, der selv foreløbigt ikke kan være Genstand for Begrundelse. \nm{18}

Denne Art Intuition (c) forekommer ikke hos Bergson, og derved forklares tildels den Modsætning, han hævder mellem Anskuelse og Forstand, en Modsætning, man ikke uden Ret har betegnet som en uheldig Begivenhed i Tænkningens Historie.

d. En fjerde Art af Intuition er den umiddelbare Skuen af en Sammenhæng eller en Helhed, som kan vindes ved opmærksom og tænkende Gennemgang af en Række eller Gruppe af Led eller Dele, med særskilt Opfattelse af deres indbyrdes Forhold. Saaledes kan Gangen i en indviklet Bevisførelse eller Sammenhængen mellem forskellige Synspunkter for et og samme Emne tilsidst staa for os som Genstand for et samlet Blik. Eller en Personlighed, hvis enkelte Egenskaber og Handlinger vi have iagttaget og overtænkt hver for sig, kan tilsidst staa for os i hele sin individuelle Ejendommelighed. Jo mere den tredie Art af Intuition (c) har været medvirkende ved den diskursive Gennemgang, des klarere og tydeligere vil Helheden senere kunne staa for os. Selve Tankearbejdet (Analyse og Bevisførelse) kommer da Helhedsanskuelsen til Gode, og de to Arter af Intuition, den tredje (c) og den fjerde (d), ere ingenlunde Modsætninger. Den ene kan bane Vejen for den anden. Det er den ad Tankens Vej fundne lovmæssige Sammenhæng, der
% --- p. 25 ---
\opage{25}betinger Indsigten i Helheden. Baade hos Descartes og hos Poincaré bruges Udtrykket Intuition da ogsaa, ikke blot om de enkelte Tankeovergange, men ogsaa, og i nøje Sammenhæng dermed, om det Overblik, ved hvilket et Tankehele tilsidst kan staa for Bevidstheden.

Spinoza gav denne Art Intuition Plads som Højdepunktet af menneskelig Erkendelse. Den er forskellig baade fra den blotte Iagttagelse af de enkelte Emner og fra den abstrakte Erkendelse af Lovene for deres Optræden. Den gaar ud paa at se det Enkelte i dets individuelle Ejendommelighed og dog som et Sammenspil, en Samfoldighed af almene Love, gennem hvilke det staar i indre Forbindelse med den øvrige Tilværelse. Det er for Spinoza Erkendelsens Fuldendelse, men han erklærer, at det kun er meget Lidet, han har kunnet forstaa paa denne Maade. \nm{19} Han har ved sin Beskrivelse af denne Art Intuition snarere stillet en stor, tilsidst uløselig Opgave, end selv givet Anvisning til dens Løsning. Og en saadan Forening af det Individuelle og det Almene, som han fordrer, ligger i hvert Tilfælde paa Grænsen mellem Videnskab og Kunst. ---

Bergson’s „metafysiske Intuition“ kan minde om denne fjerde Art af Intuition, men den betoner ikke saa bestemt som denne Sammenhængen mellem Tankevirksomhed og afsluttende Skuen. Den metafysiske Intuition skal blive til ved at uddybe og udvide (creuser et élargir) Sanseanskuelsen (la perception), uden at der angives nogen bestemt Metode for denne Uddybelse og Udvidelse. Tankearbejdets enkelte Arter og Overgange blive ikke af nogen Betydning; det rationelle Element skydes tilside --- og derfor nærmer Bergson’s metafysiske Intuition sig snarere til den anden (b) end til den fjerde (d). Paa romantisk Vis betragtes derhos Intuitionen som en Genvinden af noget Tabt, en Genvinden, der sker mere paa Trods af Forstanden end ved dens Hjælp. I Modsætning hertil maa det hævdes, at forsaavidt Intuitionen af fjerde Art (d) har videnskabelig Betydning, peger den fremad ikke tilbage. Den bliver mulig ved tankemæssig Bearbejdelse af Intuitioner af første Art (a), gennem de Analyser og Deduktioner, Intuitioner af tredie Art (c) bane Vej for. Jo rigere Intuitionen a er, og jo strengere gennemført
% --- p. 26 ---
\opage{26}Intuitionen c er, des betydningsfuldere kan Intuitionen d blive. Hvis Filosofien kunde naa en saadan Afslutning, som man i de store Systemers Tid antog, saa maatte det være gennem en paa saadan Vis vunden Intuition, at Maalet skulde kunne naas.

Medens Bergson’s psykologiske Intuition (af Arten a) er af afgjort Betydning og har ført ham til at sige et Ord i rette Tid, er det ikke lykkedes ham at give et klart Begreb om den metafysiske Intuition, hvis Mulighed han hævder, eller at begrunde dens filosofiske Gyldighed og Værdi.

\chapterend

% --- p. 27 ---
\chapterhead{PSYKOLOGI OG FYSIOLOGI}

\opage{27}13. \initial{B}{ergson’s} Filosofi, især hans Lære om Intuition, er for en stor Del rettet mod Naturvidenskabens Indflydelse paa Aandsvidenskaben. Men saa knytter der sig en særlig Interesse til den Maade, paa hvilken Forholdet mellem Fysiologi og Psykologi, eller mellem Legeme og Sjæl opfattes i denne Filosofi. Thi her træder Forholdet mellem de to store Afdelinger af menneskelig Videnskab bestemtest og skarpest frem. I sit Værk „Matière et Memoire“ % PRINTED AS IS: Memoire, no accent
(1897) drøfter han dette Spørgsmaal.

14. Bergson begynder med at stille sig paa et strengt naturvidenskabeligt Standpunkt. Han vil holde sig til, hvad objektiv Iagttagelse viser, og hvad naturvidenskabelig Metode godtgør. Alt i Naturen bestaar i Bevægelse. Hvad der sker i Omverdenen og i mit eget Legeme, Hjerne og Nervesystem indbefattede, er kun Bevægelse af forskellig Art og Grad. Fra Omverdenen forplanter Bevægelsen sig ind i vort Legeme. Den naar gennem Sanseorganer og Nervetraade til Rygmarv og Hjerne og fortsættes i udadgaaende Bevægelser. Vort Legeme er et Redskab til Modtagelse af Bevægelse udefra og til Udløsning af Bevægelse udadtil. Der er i denne Henseende ingen Forskel mellem Hjernen og de andre Dele af Legemet. Kun har Hjernen mere end de andre Organer Evne til at bevare modtagne Paavirkninger saaledes, at de ikke strax føre til udadgaaende Bevægelser, men kunne komme til at medvirke ved senere Bevægelser. Legemet, Hjernen ind%
% --- p. 28 ---
\opage{28}befattet, er et Bevægelsesredskab, ikke Andet. Ikke i nogen Grad, ikke i nogen Betydning, ikke fra noget Synspunkt tjener den til at forberede, og endnu mindre til at forklare en Forestilling (représentation). Hvad der i vore Sanseiagttagelser (perception) kan forklares ved Hjernens Virken, det er de Handlinger, som begyndes eller forberedes eller foranlediges, men det er ikke selve vore Sanseiagttagelser. Endnu mindre indeholder Hjernevirksomheden nogen Forklaring af Erindring og højere Aandsvirksomhed. Og naar Alt i og udenfor Hjernen er Bevægelse, saa maa det føjes til, at man i Bevægelsen ikke maa søge Andet, end hvad man iagttager. Nervesystemet har kun rent fysiske Egenskaber og har ikke Evne til Andet end at modtage, hæmme eller fortsætte Bevægelse. Vi sanse kun Bevægelse --- og Bevægelse kan ikke frembringe Andet end Bevægelse. De materielle Processers hele Virkning er udtømt i Arbejdet paa Tillempelse af Bevægelser (travail d’adaptation motrîce). % PRINTED AS IS: motrîce (motrice)
Beskadiges Hjernen, saa er det kun Bevægelser, der hæmmes eller hindres, ikke Andet \nm{20}.

Saadanne Sætninger kunde man vente at finde hos en konsekvent Materialist. Men Bergson adskiller sig fra Materialismen ved at nægte, at materielle Bevægelser kunne bevirke Sjæle- eller Bevidsthedsliv. Det Spørgsmaal melder sig da med fornyet Eftertryk, hvorledes efter hans Opfattelse Bevægelse i Rummet, særligt i Hjernens Smaadele, forholder sig til de sjælelige Tilstande.

15. At fornemme og at genkende betyder for Bergson i simpleste Tilfælde kun at mærke en Forskel mellem det, der har Betydning for vor Trang, og det, der ikke har det, og at man mærker denne Forskel, viser sig i, at den bestemmer Bevægelserne, som udløses. Men dette fører os, mener han, ikke ud over det Materielle. Fornemmelse er kun en Del af Materien, eller der er dog kun en Gradsforskel mellem Fornemmelsen af Materien og selve Materien. Det er ligegyldigt, om vi her tale om Fornemmelse (sensation), Iagttagelse (perception) eller Genkendelse (reconnaissance). Der gives, som Grænsetilfælde, en øjeblikkelig Genkendelse, som Legemet paa egen Haand
% --- p. 29 ---
\opage{29}(tout seul) er istand til, uden at nogen egentlig Erindring kommer til; den bestaar i Handling, ikke i Forestilling (représentation).

Vort i Rummet udstrakte Legeme har altsaa baade Evne til Fornemmelse og til Bevægelse. Den materielle Side af vor Existens er en Helhed af Fornemmelse og Bevægelse \nm{21}. ---

Men --- spørger man uvilkaarligt --- naar Fornemmelse (Iagttagelse, Genkendelse) kun er Bevægelse, ligesom alt Andet i „den materielle Side af vor Existens“ (la matérialité de notre existence), hvorfor bruger saa Bergson to Betegnelser, hvorfor skelner han overhovedet mellem Fornemmelse og Bevægelse? Han synes her at glemme sin egen Advarsel mod Sproget, der i dette Tilfælde maatte have dannet et Ord for meget. At der foruden Bevægelse ogsaa existerer Fornemmelser, kan han ikke vide af ydre Iagttagelse, der, som han med Rette bemærker, kun viser os Bevægelser i Rummet. Her maa Selviagttagelse træde til. Det er i Kraft af den, at vi f. Ex. skelne mellem Smerten og de Vridninger af Kroppen, hvori den ytrer sig. Kunde vi iagttage Hjernen i alle dens forskellige Tilstande, vilde vi kun se Bevægelser og Omlejringer af Atomer, ikke Andet. At der samtidigt føles eller tænkes, maa vi vide ad anden Vej.

Naar Bergson uden videre tillægger Legemet Evne til at skelne (fornemme, iagttage, genkende), saa tager han i Virkeligheden sin første Sætning om, at Alt i den materielle Verden bestaar i Bevægelse, og at Bevægelse kun kan frembringe Bevægelse, tilbage. Eller med andre Ord: indenfor det, han kalder Skelnen eller Fornemmelse, dukker hele det store Problem om Forholdet mellem Sjæl og Legeme op igen i Form af det Spørgsmaal: hvorledes forholder Fornemmelsen i Betydning af Bevægelse sig til Fornemmelse som Bevidsthedsfænomen (la perception en tant que consciente)? Og han maa da indrømme, at hvad Hjernen kan forklare, „hvad vore Fornemmelser angaar, det er kun vore begyndte eller forberedte Hindringer, % PRINTED AS IS: Hindringer for Handlinger (Bergson: nos actions commencées ou préparées; cf. p. 28)
ikke selve vore Fornemmelser \nm{22}.“

Saaledes gaar det ofte, naar man vil omgaa de store Problemer. Selv om de reduceres til tilsyneladende smaa Dimensioner, bevare de dog deres Brod.

16. Medens der ifølge Bergson kun skal være en Grads%
% --- p. 30 ---
\opage{30}forskel mellem Fornemmelse og (anden) Materie, hævder han en Artsforskel mellem Fornemmelse og Erindring.

Det er Erindringen, der for Bergson godtgør Kontinuiteten, den indre Sammenhæng i Sjælelivet. Vi kunne ofte iagttage, hvorledes Erindringsbilleder voxe frem og tiltage i Klarhed og Tydelighed. Her staa vi overfor den indre Strøm, som Bergson saa ofte henpeger paa, og som han formaar at beskrive paa saa fængslende en Maade. Naar der ofte synes os at være større Sammenhæng mellem de ydre Genstande end mellem vore indre Tilstande, saa paastaar Bergson, at det beror paa unøjagtig Iagttagelse. „Ved at se nærmere til vil man opdage, at vore Erindringer danne en Kæde, og at vor Karakter, som stedse er tilstede i vore Beslutninger, er den aktuelle Syntese af alle vore tidligere Tilstande. I denne fortættede Form har vort forbigangne sjælelige Liv endog en højere Grad af Existens for os end den ydre Verden; af denne opfatte vi stedse kun en lille Del, medens vi gøre Brug af hele Indbegrebet af vor gennemlevede Erfaring.“ \nm{23}

Istedetfor den sædvanlige Modsætning mellem Materie og Sjæl stiller Bergson altsaa Modsætningen mellem Fornemmelse og Erindring.

Men hvad hjælper det, naar „Fornemmelsen“ ifølge Bergson i sin simpleste Form er rent materiel? Han kalder jo ogsaa sin Bog „Materie og Erindring“, og med Rette, naar „Fornemmelsen“ hører ind under det Materielle. Spørgsmaalet bliver stadigt det gamle: hvorledes kan Bevidsthed (Erindring) opstaa af materiel Bevægelse eller gaa over til materiel Bevægelse? ---

Og i Grunden nægter Bergson selv, at der gives rene, det vil sige, absolute, isolerede, selvstændige Fornemmelser. Kun som et Grænsebegreb, et idealt Tilfælde skal man kunne tale om saadanne. De kunde kun opfylde et enkelt Øjeblik; men Bergson nægter jo, at Tiden kan opløses i absolut enkelte Øjeblikke. Den nyere Psykologi har, i Fechners Fodspor, godtgjort, at enhver Fornemmelse er bestemt, ikke blot ved den nærværende Sansepaavirkning, men ogsaa ved de umiddelbart forudgaaende. Allerede her virker Fortiden med til at bestemme
% --- p. 31 ---
\opage{31}Nutiden, saa at man med Rette kan kalde enhver Fornemmelse for en elementær Erindring. Det er ikke blot i vore Beslutninger, saaledes som Bergson træffende har beskrevet, men det er i alle vore Oplevelser, selv de mest elementære, at Fortiden virker med. Bergson’s Kritik af Fechner har ført ham for vidt. Selv om det ikke lykkedes Fechner at grundlægge en matematisk Psykologi, saa vise de af ham fremdragne Kendsgerninger dog, at der ogsaa udenfor den egentlige Erindring er en direkte Sammenslutning mellem vore Oplevelser. Bergson havde altsaa kunnet udstrække sin Opfattelse af Erindringen som Udtryk for Sjælelivets Sammenhæng længere ned, til de mest elementære sjælelige Funktioner. Og han indrømmer faktisk, at i hine virkelige, konkrete Fornemmelser medvirker Erindringen: dans la perception concrète la mémoire intervient! --- Den „konkrete Sanseiagttagelse“ sættes udtrykkeligt i Modsætning til den rene, ideale Fornemmelse, som vilde finde Sted i et absolut Øjeblik, men som existerer „en droit plutôt qu’en fait!“ \nm{24}

Artsforskellen mellem Fornemmelse (Sanseiagttagelse) og Erindring maa derfor opgives. At der overhovedet kan gøres en Forskel mellem Fornemmelse (Iagttagelse) og Erindring, skyldes Erfaring gennem de Skuffelser, som belære om, at man ikke kan handle efter alle de Forestillinger, der dukke op. Enhver Erindring har fra først af, ligesom Fornemmelsen, Tendensen til at udløse Handling. Erindringens første Form er Forventning eller Drift. Vi begynde med Tro til alle vore Forestillinger, saavel som til alle vore Fornemmelser. Det er i Erfaringens Skole, vi lære at henvise nogle Forestillinger til Fortidens Verden en Gang for alle, medens andre Forestillinger fremdeles bestemme vore Forventninger og vore Formaal. \nm{25}

17. Naar Bergson søger at fastholde Kontinuiteten saavel indenfor Sjælelivet, som indenfor den materielle Natur, hist ved Hjælp af Begrebet Erindring, her ved Hjælp af Begrebet Bevægelse, kunde det have ligget nær for ham at slutte sig til den af Spinoza grundlagte Hypotese, der antager et Forhold af Proportionalitet eller Analogi mellem Bevidsthed og Bevægelse,
% --- p. 32 ---
\opage{32}uden at dog Bevidsthed skal kunne udledes af Bevægelse eller Bevægelse af Bevidsthed. Det er en Hypotese, som ikke giver nogen Løsning, men som angiver tydeligt, hvorledes vi maa opfatte Forholdet, før vi indlade os paa videregaaende Spekulationer. Bergson anerkender den som en Arbejdshypotese, hvad der i Virkeligheden er en stor Anerkendelse, da en Hypoteses Berettigelse især maa vise sig i den Brug, man kan gøre af den. Og desuden er det tilsidst en Umulighed, at en „Arbejdshypotese“, der gennemgaaende viser sig anvendelig, idet den stiller nye Opgaver, og idet de fundne Resultater stemme med den, ikke skulde have mere Betydning end et Redskab, man lader falde, naar man ikke har Brug for det. Det er et Bidrag til Karakteristik af Tilværelsen, at vi kun kunne forstaa den ad aldeles bestemte Veje og ud fra visse bestemte Antagelser. En endnu større Anerkendelse yder Bergson den nævnte Hypotese, naar han indvender imod den, at den kun er „en forklædt Konstateren af Kendsgerningerne“. Hvori Forklædningen egentligt bestaar, véd jeg ikke; men en Fremstilling af Kendsgerningerne er dog ikke saa ilde et Grundlag for videregaaende Eftertanke \nm{26}.

At Bergson ikke lægger større Vægt paa saadan „Konstateren“, hænger sammen med, at Naturvidenskaben for ham er en Træl i det praktiske Livs Tjeneste og derfor ikke har Krav paa at yde principielle Bidrag til vor hele Verdensanskuelse. Og dog er det klart, at Problemet om Forholdet mellem det Aandelige og det Materielle først faar sin egentlige Brod ved den Kendsgerning, at Naturvidenskaben har udviklet sig til en Lære om den materielle Natur som en stor, kontinuerlig og uoverskuelig Sammenhæng af Bevægelsesprocesser, --- hvad jo Bergson ogsaa selv klart har set. Han har endog bidraget Meget til at skærpe den Indsigt, at man ikke har Ret til at antage, at hvad der foregaar i Hjernen er Andet end Bevægelser, ikke Processer af mere eller mindre mystisk Art. Men deraf drager han saa den uberettigede Slutning, at Hjernen og Organismen slet ingen Betydning have for Bevidsthedslivet uden netop i de Øjeblikkke, % PRINTED AS IS: Øjeblikkke (three k)
hvor det skal gribe ind i den materielle Verden: ellers fører Bevidstheden sit Liv paa egen Haand, og
% --- p. 33 ---
\opage{33}et Liv, der omfatter langt flere Muligheder, end den materielle Natur kan realisere.

I en nyligt udkommen Afhandling om Sjæl og Legeme har Bergson herom udtalt sig saaledes: „Den, der kunde se ind i Hjernens Indre, naar den er i fuld Virksomhed, følge Atomernes Kommen og Gaaen og tyde Alt, hvad de gøre, vilde ganske vist kende Noget af det, der foregaar i Sjælen; men det vilde kun være meget Lidet, han kendte. Han vilde kende netop det, som kan udtrykkes i Legemets Miner, Stillinger og Bevægelser, --- det i en sjælelig Tilstand, som bestaar i en Handling, der er ved at fuldbyrdes eller maaske blot ved at fødes; alt Andet vilde undgaa ham. Han vilde overfor de Tanker og Følelser, der udfolde sig i Bevidsthedens Indre, være i samme Situation som en Tilskuer, der tydeligt ser Alt, hvad Skuespillerne gøre paa Scenen, men ikke forstaar et Ord af, hvad de sige.“

Ved denne Opfattelse krænkes egentligt baade Psykologien og Fysiologien. --- Psykologisk set kan Villie og Handling ikke isoleres fra de andre Sider ved Sjælelivet. En Handling „fødes“ der i hvert eneste Øjeblik; ti der er ikke en Fornemmelse, ikke en Følelse, som ikke medfører et Villieselement. Det er en falsk Abstraktion, naar Erindringens Liv sondres fra Handlingen paa den Maade, Bergson gør det. Han gør Bevidsthedslivet altfor meget til et Drømmeliv. Og selv i Drømme, selv i de „fromme“ Ønsker, hvortil den oprindelige sangvinske Handletrang kan omsættes, er den stadige Handlingstilskyndelse ikke helt forsvunden. Sjælen er her, for at bruge et Billede af Helen Keller, som en Fugl, der vil arbejde sig op imod Vinden. Naar Handlingens Øjeblikke have fysiologiske Paralleler, saa maa derfor ogsaa alle andre sjælelige Tilstande have det. Og det maatte da nok for en psykologisk kyndig Iagttager være muligt at tyde Hjærneprocesserne % PRINTED AS IS: Hjærne- (Hjerne-, as elsewhere in the book)
ad Analogiens Vej, selv om maaske de Øjeblikke, hvor der skrides til udadgaaende Handling, vilde være de lettest forstaaelige. --- Fysiologisk set maa der overfor den Paastand, at der i Sjælen indeholdes langt Mere, end Legemet kan udtrykke, mindes om en Sætning af Spinoza, hvis Gyldighed stadigt staar ved Magt:
% --- p. 34 ---
\opage{34}„Ingen har endnu af Erfaring lært, hvad Legemet, betragtet som blot Legeme, kan virke efter Naturens egne Love, og hvad det ikke kan. Ti Ingen kender endnu Legemets Indretning saa nøjagtigt, at han skulde kunne forklare alle dets Funktioner.“ --- Det er netop den store Betydning af Spinoza’s Opfattelse paa dette Punkt, at den paa én Gang opfordrer baade Fysiologien og Psykologien til fortsat Arbejde paa at efterspore den indre Sammenhæng paa deres to Omraader, for derigennem langsomt at bane Vej for en Forstaaelse af disse to Omraaders indbyrdes Forhold. ---

I den Sætning, at Sjælelivet indeholder Mere, end Legemet nogensinde kan udtrykke, Legemet, der kun er et Redskab for Sjælen, naar den skal handle, finder Bergson Muligheden for Sjælens Udødelighed; men det Spørgsmaal, om denne Udødelighed gælder for en uendelig eller kun for en begrænset Tid, en Tid x, kan ikke besvares ad videnskabelig Vej. Hans Tankegang her falder sammen med en Hypotese, som Kant, (skønt den kun var „ein bleyernes Waffen“) anbefalede til polemisk Brug overfor dogmatiske Materialister, og som William James senere har optaget. Legemet opfattes da, paa platonisk Vis, som en snever Kanal, i hvilken Sjælelivets Strøm presses ind, og som stadigt begrænser denne, uden at kunne stanse den og uden at kunne hindre den i at forene sig igen med Havets store Strømme. \nm{27}

\chapterend

% --- p. 35 ---
\chapterhead{UDVIKLINGSFILOSOFI}

\opage{35}18. \initial{D}{et} er et Tidernes Tegn, at enhver Filosof i vore Dage maa tage Stilling overfor Udviklingsbegrebet. Dette Begreb har nu indtaget sin Plads mellem Tankens Grundformer eller Kategorier, idet det indeholder Anvisning til at stille Spørgsmaal og Opgaver. Overfor enhver Begivenhed og ethvert Fænomen maa der spørges, gennem hvilke Stadier de ere naaede frem til deres foreliggende Beskaffenhed. Det er en speciellere Form for det almindelige Aarsagsbegreb. \nm{28} En Filosofi karakteriseres fra en væsentlig Side ved den Stilling, den giver dette Begreb, og ved den Maade, hvorpaa den anvender det.

For Bergson maatte det særligt ligge nær at give dette Begreb en fremragende Plads. Han hævder jo, at baade det umiddelbart Givne (det, vi staa overfor og selv ere Del af, naar vi arbejde os ud af den udstykkende, mekaniske Betragtningsmaade, Livets Praxis paanøder os) og den afsluttende Anskuelse, vi naa til gennem Tankearbejdets Skærsild, viser os en kontinuerlig Bevægelse, en fremskridende Strøm, indenfor hvilken nye Egenskaber og Ejendommeligheder opstaa og udfolde sig. Alt er for ham Virksomhed og Udvikling; der gives ingen død Væren, ingen faste Ting, ingen absolut Hvile.

Men den naturvidenskabelige Opfattelse af Udviklingsbegrebet er ifølge Bergson ufyldestgørende. Den lægger altfor stor Vægt paa de ydre Forhold. Udviklingen kommer efter den til at staa som en Sammenstykning af Elementer, der forud bestod som isolerede, --- som en Forbindelse af enkeltvis fremkomne Tilløb. En virkelig, indre Sammenhæng bliver der her ikke Tale om.
% --- p. 36 ---
\opage{36}Denne udvortes Betragtningsmaade er for Bergson en Følge af Naturvidenskabens Oprindelse og stadige Opgave. Den er udsprungen af Livets Behov, og dens Opgave er at tjene Livets Bestaaen og Udvikling. Men dens Opgave er ikke at forstaa selve Livet, i hvis Tjeneste den staar. Den er frembragt af Livet under visse bestemte Forhold og for at virke paa visse bestemte Ting; hvorledes skulde den da kunne omfatte Livet selv, af hvilket den kun er en enkelt Ytring, og for hvilket den kun afgiver et enkelt Synspunkt? Den er aflejret undervejs i Løbet af Udviklingsprocessen; hvorledes skulde den da kunne finde Anvendelse paa selve denne Proces i dens Helhed?

Spencer benyttede den Kendsgerning, at Videnskaben er opstaaet under Udviklingen gennem Kamp for Tilværelsen i utallige Generationer som Bevis for Gyldigheden af de Begreber, der ligge til Grund ved det videnskabelige Arbejde. Men Bergson finder i denne Kendsgerning netop en Indvending mod disse Begrebers absolute Gyldighed. Der er, mener han, ingen Grund til at tro, at den Opfattelsesmaade, som i Udviklingens Løb er paatvunget os af praktiske Grunde, under Indflydelse af Livets Trang og Krav, skulde være den eneste mulige. Maaske ere heller ikke alle aandelige Kræfter tagne i Beslag af det praktiske Behov. Til de ubrugte Kræfter maa man da vende sig, hvis man vil naa en højere, uinteresseret Forstaaelse af Livet og dets Love! Her ligger ifølge Bergson Filosofiens Opgave. \nm{29}

Men før han forsøger at løse denne Opgave, drøfter han to Hovedformer, i hvilke den udvortes og mekaniske Opfattelse af Livets Udvikling, han bekæmper, ere fremtraadte.

19. Livet bestaar under stadig Kamp med Materien. Det gælder om at behandle materielle Ting, bemægtige sig dem og omforme dem. For at leve behøver man kun Evne til en saadan Behandling af Materien; derimod er det ikke nødvendigt for at leve, at man forstaar, hvad Livet selv egentligt er. Det er da intet Under, at vor Forstand kommer tilkort her. Det rent Materielle kan opfattes som sammensat af Dele, ved hvis Sammenføjning det bliver til. Ved en rent materiel Helhed have vi
% --- p. 37 ---
\opage{37}Delene givne først, og Helheden er et simpelt Resultat af deres Sammenføjning.

Naar nu Videnskaben staar overfor et levende Væsen, er det naturligt, at den prøver at anvende den samme Betragtningsmaade overfor det, som den er vant til at anvende overfor alt rent Materielt. Ogsaa Livet skal da være et blot Produkt, opstaaet ved Sammenføjning af Dele. Man søger da at reducere Læren om Livet til Fysik og Kemi. Det er jo naturligt, naar det blot drejer sig om at øve en saadan Indvirkning paa det Levende, at det kan tjene som Middel for bestemte Formaal. Man opfatter al Organisation som en Slags Fabrikation, og enhver Organisme som en Maskine, hvis Dele og hvis Bygningsmaade det gælder at finde. Voilà le point de vue de la science! --- Det er i Konsekvens af dette Synspunkt, at mange Videnskabsmænd se det som Idealet at underlægge Organismen samme matematiske Behandling, som Solsystemet forlængst er bleven underlagt i Astronomien. \nm{30} ---

Saaledes skildrer Bergson den mekanistiske Opfattelse af Livet. Den fremgaar efter hans Mening ligefrem af den Metafysik, som Menneskene naturligt gaa ud fra, uden at være sig bevidst, at det er Metafysik. Overfor den stiller han en Række af Betragtninger, som føre i helt anden Retning, idet de skulle godtgøre, at Livet ikke kan reduceres til et System af fysiske og kemiske Processer.

Den mekanistiske Opfattelse kan ikke forklare Livets Kontinuitet, den nøje Sammenhæng mellem dets Fortid og dets Nutid. De senere Stadier i Livsudviklingen skyldes ikke Kombination af Elementer, der allerede vare til forud og nu først forbindes, men de kunne bringe noget helt Nyt. Nye Kvaliteter kunne opstaa, der ere betingede ved Fortiden, men ikke kunne udledes af denne. Ved den inderlige Sammenhæng mellem det Gamle og det Nye minder det organiske Liv om Sjælelivet. Begge Steder gælder det, at de forskellige Stadier i Udviklingen ere inkommensurable indbyrdes, og at Stadiernes Række ikke kan vendes om. En Organismes Væxt, og selve Arternes Udvikling i Naturen, frembyder en bestemt Rækkefølge, der ikke ligesaa godt kunde tænkes omvendt. Vi komme her atter til%
% --- p. 38 ---
\opage{38}bage til Forskellen mellem det psykologiske Tidsbegreb (durée vraie), ifølge hvilket alle Øjeblikke ere kvalitativt forskellige og ikke kunne byttes om, og det mekaniske Tidsbegreb, ifølge hvilket Tiden bestaar af ens Øjeblikke, der uden videre kunne sættes istedetfor hverandre, kun forskellige ved deres Nummer i Rækken. ---

Ligesaa uholdbar som den mekanistiske Teori er ifølge Bergson den finalistiske eller teleologiske, der antager en forudfattet Plan, der allerede indeholder Alt, hvad der senere kommer frem i Erfaringen under Udviklingens Forløb. Ogsaa her bliver Tiden tilsidst unyttig. Der opstaar intet Nyt. Hvad Erfaringen frembyder, bliver kun Gentagelse af, hvad der allerede var indeholdt i Planen. Finalismen er kun en omvendt Mekanicisme.

Ogsaa denne Opfattelse er opstaaet under Indflydelse af praktiske Forestillinger. Livets Krav faar os til at lægge Planer, som vi bagefter søge at sætte ind i Virkeligheden. Og vi maa da her holde os til, hvad vi kunne forudse, og det er det, der kan gentage sig, naar samme Forhold kommer igen. Men Livet er ikke udtømt med det, der gentager sig; nye Egenskaber og Former kunne opstaa, der ikke kunne forudses, skønt de maaske bagefter kunne forklares. Ja, Gentagelse er overhovedet kun mulig for en rent abstrakt Betragtning. I den virkelige Tid (durée réelle), --- den, hvis Tand efterlader Mærket af sit Bid paa Alt, --- finder ingen Gentagelse Sted. Alt forandres helt igennem, og den samme konkrete Virkelighed gentager sig aldrig.

20. I denne Kritik af de to modstaaende Teorier har Bergson en Forgænger i Kant’s berømte Tankegang i „Kritik der Urtheilskraft“, som han mærkeligt nok ikke nævner, skønt den har præciseret Problemet paa den Maade, paa hvilken det i Grunden fremstiller sig den Dag idag. Kant viste, at der for den menneskelige Erkendelse kun er to Maader at forklare en Helhed paa, nemlig enten at opfatte den som Produkt af forud bestaaende Dele, eller at opfatte den som et efter en forud given Plan frembragt Værk, --- men at ingen af disse to Maader kunde anvendes paa den Helhed, en Organisme er. Medens Kant heraf kun sluttede, at vi ikke maa forvexle vor Maade at
% --- p. 39 ---
\opage{39}erkende paa med Naturens Maade at virke paa, mener Bergson at kunne give en positiv Løsning af Spørgsmaalet. Det er især paa dette Punkt, at han forsøger at naa ud over, hvad den af Kant grundlagte Kriticisme kan yde.

Forstanden, Erkendelsen --- saaledes er hans Tankegang --- er kun et Udsnit, en Indsnævring indenfor en mere omfattende Helhed. Livet gaar ud over Forstandens Rækkevidde. Ved sin praktiske Synsmaade har Forstanden udelukket sig fra Livets store Sammenhæng. Og det er til denne store Sammenhæng, vi maa gaa tilbage. Men hvorledes kunne vi naa den? Kun ved at gennemleve den! Vi mærke, at der udenom den klare Kærne, Forstanden formaar at finde, er en ubestemt Frynse, som er ved at tabe sig i Mørket, og som ikke er til nogen Nytte i Livets Praxis, men som antyder en større Dybde, end Mekanicismen og Finalismen kende. \nm{31}

Denne vage Intuition vinder nu ifølge Bergson i Klarhed ved en nærmere Betragtning af det organiske Livs Udviklingsgang.

Under den store Brydningskamp med Materien udfolder Livet sig i forskellige Retninger og løber ud i afvigende Typer, der staa som Endepunkter for disse Retninger. Saaledes fremtræder Modsætningen mellem Plante- og Dyretypen: hin fæstnet til Sted og uden Bevidsthed, denne bevægelig og med Evne til Bevidsthed. De have udviklet sig af fælles Former. Indenfor Dyreverdenen er der igen forskellige Udviklingsrækker, af hvilke de mærkeligste er den, der udfolder sig i Insekttypen, og den, der gennem Hvirveldyrsrækken fører op til Mennesketypen. Af disse to Typer er den ene karakteriseret ved Instinktets, den anden ved Forstandens Udvikling. Instinktet er en simpel Fortsættelse af Væxten, en Evne til at gøre Brug af Organerne (de organiske Redskaber) paa en bestemt Maade; Forstanden er Evne til at frembringe mekaniske Redskaber. Mekanisk Opfindelse er Forstandens væsentlige Funktion. Den Dag idag er det sociale Liv væsentligt bestemt ved Fabrikation og Udnytten af kunstige Redskaber, og de Opfindelser, der som Milepæle betegne Fremskridtets Vej, have ogsaa betinget dets Retning. Instinkt og Forstand ere to forskellige, men lige elegante Løsninger af et og samme Problem.
% --- p. 40 ---
\opage{40}I Plantelivets Vegeteren, i Instinktets uvilkaarlige Virken og i Forstandens bevidste Fabrikeren har Livsprocessen udfoldet sig ad tre forskellige Veje. Ligefra Aristoteles’ Tid af har man atter og atter begaaet den Fejl at betragte disse tre Livstyper som tre successive Trin indenfor en og samme Udviklingsrække. Men der er mellem dem ikke en Gradsforskel, men en Artsforskel. Det er tre indbyrdes afvigende Retninger for en Virken, der maatte spalte sig for at kunne voxe. \nm{32}

Det er overhovedet en Ejendommelighed for Livet, at det ad meget forskellige Veje formaar at naa samme Resultat. Saaledes muliggøres Forplantning paa meget forskellig Maade indenfor de forskellige Typer. Synsevnen frembyder især et tydeligt Exempel paa den nævnte Ejendommelighed, naar man sammenligner de forskellige Udviklingsformer af Øjet hos forskellige Dyr. \nm{33}

Alt dette vidner om en Stræben, en Higen, en oprindelig Livsimpuls eller et Livsopsving (élan de vie), der baner sig Vej gennem Materiens Verden, og som i de mangfoldige Typer, saavel som i de enkelte Individer hvert for sig, afsætter de Resultater, den har vundet trods de Skavanker, den stedse brydes imod. Livets Historie er Udfoldelse af en oprindelig Kraftbeholdning, en Udstrømmen af en Kilde, hvis Vande søge at uddybe sig Flodlejer i Materiens Masser. Materien i sig selv har, som den saakaldte Entropi viser, en stadig Tendens til at henfalde i absolut Ligevægt. Det er imod denne Tendens, Livet stadigt kæmper. Livsudviklingen er ikke skabende i absolut Betydning; den er stedse betinget ved Forholdet til de materielle Kræfter og Masser. Men der ytrer sig i den Noget, som hverken Mekanicismen eller Teleologien lader komme til sin fulde Ret, --- en frugtbar Enhed, en uendelig Fylde. \nm{34} ---

21. Som man ser, gaar Bergson ud fra en absolut Modsætning mellem det Organiske og det Uorganiske. Saa store Problemer der end her stadigt gøre sig gældende, er det dog et stort Spørgsmaal, om en Filosof gør vel i at tage Udgangspunktet for sin hele Anskuelse i en saadan Modsætning. En Kendsgerning er det, at hvor meget Videnskaben end stadigt
% --- p. 41 ---
\opage{41}maa fastholde Forskellen mellem Livet og det Uorganiske, saa betinger dette dog ikke nogen Forskel i den videnskabelige Metode. Man anvender overfor det Levende samme Metoder som overfor det Livløse, og kun forsaavidt man gør det, naar man holdbare Resultater. Og det er en ikke ualmindelig Opfattelse, der f. Ex. er udtalt af Jacques Loeb \nm{35}, at „kun tekniske Vanskeligheder ere Skyld i, at kunstig Frembringelse af levende Materie endnu ikke er lykkedes.“

Men selv om man indrømmer den skarpe Modsætning mellem Organisk og Uorganisk, er Bergson’s Opfattelse ikke strengere begrundet end de to Opfattelser, han afviste. Forholdet er dette, at der overfor et Emne som det organiske Liv stedse maa opereres med Analogier. Det er en Analogi, Mekanicismen anvender (saalænge „de tekniske Vanskeligheder“ ikke ere overvundne), naar den opfatter Organismen som en Maskine. Det er en Analogi, Teleologien anvender, naar den forklarer Livsfænomener som Frugt af Hensigt og Plan. Men det er ogsaa en Analogi, Bergson bruger, naar han opfatter Livet som „Opsving“, Stræben, Higen. Han adskiller sig kun fra Teleologien (med hvilken han desuden har en vis Sympati) ved at hente sin Analogi fra det uvilkaarlige, halvt ubevidste Sjæleliv istedetfor fra det klart bevidste, forstandigt beregnende Sjæleliv. Livet er for ham egentligt af psykologisk Natur. Hans Lære er en Art Vitalisme, bygget paa Antagelsen af en særegen Livskraft, der brydes med de fysiske Naturkræfter, en Antagelse, der mere er et Udtryk for Beundring og Forundring overfor Livets mægtige og rige Udfoldelse end en Anvisning til nøjagtig Forklaring af Livets specielle Fænomener. Og naar han tyder denne Livskraft som beslægtet med det Sjæleliv, vi kunne fordybe os i hos os selv \nm{36}, saa tyder han, ligesom den romantiske Naturfilosofi, Naturen ovenfra nedad. Livet er en Kamp med Materien, en Kamp, der, forsaavidt den lykkes, fører trinvis opad mod Aandens Verden. Det er Poesi, ikke Videnskab, og skal i Grunden efter Bergson heller ikke være dette. Det er paa en dristig Analogi, han her bygger, ligesom --- efter hvad der senere vil vise sig --- hans Verdensanskuelse idethele bygger paa en Analogi, der ikke kan finde Erfaringsbekræftelse.
% --- p. 42 ---
\opage{42}Fra Schelling’s romantiske Naturfilosofi adskiller Bergson’s Naturopfattelse sig navnligt paa to Punkter. For det Første overlader han den uorganiske Natur ganske til den mekaniske Opfattelse. Den er rent materiel, og dens Hovedlov er den Carnot’ske Sætning om den stadige Tendens til Fordeling af Energien og derved til Overgang i Ligevægt. Det er med denne Tendens, at Livets Stræben stadigt har at kæmpe. For det Andet antager han en real Udvikling gennem Tiden, medens Romantikens Udviklingstanke var helt formel eller ideal, egentligt kun betød en systematisk Sammenhæng mellem Naturformerne, ikke en i Tiden fremskridende Opstaaen af de forskellige Former under visse bestemte Betingelser. \nm{37} I Mellemtiden mellem Schelling og Bergson har Udviklingsbegrebet faaet sin realistiske Karakter, der ikke mere kan fornægtes. ---

I sine mest karakteristiske Udviklingsrækker har Livsprocessen ifølge Bergson ført dels til det rene Instinkt, dels til den rene Forstand. Men for ham er Livet ikke udtømt i disse Former. Det er muligt at leve sig ind i den dybe Sammenhæng, af hvilken baade det blinde Instinkt og den analyserende Forstand kun er Udsnit. Naar vi opgive de praktiske Motiver, der have ledet Udviklingen, og pleje den uinteresserede Sans for Livets rige Virkelighed, da bliver den Intuition mulig, der for Bergson er Højdepunktet i Sjælelivet, og som jeg allerede foreløbigt har omtalt.

Før jeg anstiller en afsluttende Drøftelse af Bergson’s Intuitionsbegreb, indskyder jeg en Undersøgelse af to speciellere psykologiske Spørgsmaal --- om Villien og om Latteren --- der hvert paa sin Maade ville belyse hans Hovedtanker.

\chapterend

% --- p. 43 ---
\chapterhead{VILLIENS OG LATTERENS PSYKOLOGI.}

\opage{43}22. \initial{B}{ergson’s} Hovedbetydning ligger paa den beskrivende Psykologis Omraade. Han søger med den store Fremstillingskunst, han raader over, at føre os ud over de Distinktioner, praktiske og teoretiske Grunde føre til at gøre, ud over al Analyse, der sønderriver Sjælelivets Væv, tilbage til vort indre Livs udelte Sammenhæng, og dermed til den Kilde, hvoraf vi, uden altid at være os det bevidst, stadigt øse, naar vi tænke over psykologiske Spørgsmaal. Hvert enkelt Bæger, der fyldes fra denne Kilde, rummer kun en lille Del af, hvad der vælder frem, og medens vi øse, gaar Strømmen stedse videre. Det vilde være umuligt at genfrembringe hele Strømmen ud af disse Bægres Indhold, og tillige viser det sig, at to Bægres Indhold aldrig er ens; de ere jo tagne til forskellige Tider og paa forskellige Stadier af Strømmens Løb. Dette Billede udtrykker, skønt det ikke er hentet fra Bergson selv, et Hovedpunkt i hans Filosofi, Paastanden om en dybt liggende Kontinuitet i Sjælelivet, en Kontinuitet, der ikke maa forvexles med Enshed, da der stadigt kommer noget Nyt frem, skønt i umiddelbar Sammenhæng med og i indre Forhold til det Foregaaende.

Bergson’s psykologiske Opfattelse er, som det ses heraf, rettet mod to Sider, dels mod Udstykning af Sjælelivet i isolerede, mekanisk sondrede Elementer, dels mod Gentagelsens og Vanens Indflydelser. Han nægter ikke, at der kan ske Isolationer og Opløsninger i Sjælelivet, heller ikke, at der kan danne sig Skorper eller faste Aflejringer. Men han paastaar, at
% --- p. 44 ---
\opage{44}begge Dele kun ere Overfladefænomener, og at der, saalænge Livet varer, er en stadig, sammenhængende Bevægelse længere inde.

I to specielle psykologiske Undersøgelser fremtræder Bergson’s almindelige Opfattelse paa særligt karakteristisk Maade, og det vil derfor være af Interesse at fremdrage dem.

23. Begrebet Villie anvendes i den moderne Psykologi i noget forskellig Betydning. For Nogle er det en mystisk Kraft, ligesom „Livskraften“ i Fysiologien, og de anse det for rettest at udelukke det fra Psykologien. For Andre betegner det visse meget sammensatte og afledede sjælelige Fænomener, idet Villie efter deres Sprogbrug kun bruges om Forsætter eller Beslutninger, der forudsætte Overvejelse af forskellige Muligheder; det har heller ikke efter denne Opfattelse en berettiget Plads blandt de psykiske Elementer. Ifølge en tredie Opfattelse, som jeg for mit Vedkommende hylder \nm{38}, omfatter Begrebet Villie den hele Række af Fænomener, der begynder med Sjælelivets allerførste Ytringer og gennem Reflexbevægelser og Instinkthandlinger fører op til Drifter, Forsætter og Beslutninger. Det direkte eller indirekte Forhold til Handling af en eller anden Art og Grad mangler efter denne Opfattelse ikke i nogen sjælelig Tilstand, og alle Fornemmelser, Forestillinger og Følelser faa deres særegne Præg i de enkelte Tilfælde gennem dette Forhold, og der er derfor et Villieselement tilstede i enhver sjælelig Tilstand. --- Bergson’s Opfattelse falder nærmest sammen med denne tredie Synsmaade, kun at den fordunkles noget hos ham ved den skarpe Modsætning, i hvilken han stiller Erindringen, der for ham er det egentlige Sjæleliv, til Handlingen, der bringer Sjælelivet i Berøring med Materien.

Det er Spørgsmaalet om Villiens „Frihed“, der fører Bergson ind paa Villiens Psykologi. Han hævder, at enhver Afgørelse eller Beslutning, der med Rette skal kaldes fri, i Modsætning til, hvad der skyldes overvældende og tvingende Indflydelser og Paavirkninger eller rent øjeblikkelige Indskydelser, maa staa i inderlig Sammenhæng med hele Sjælelivet. Saa forskellige de enkelte Villiesretninger end kunne være indbyrdes, saa ere de dog Led i en og samme Udviklings%
% --- p. 45 ---
\opage{45}gang, forskellige Stadier i én Historie. Vore Oplevelser ere jo overhovedet kontinuerlige, gaa gennem mange Overgange over i hverandre, og det er rent kunstige Forskelle, vi gøre, naar vi for Bekvemmeligheds Skyld sondre stærkt imellem dem. En Personligheds Historie kender man ikke ved at holde sig til ydre Vaner eller løsrevne Ytringer. Den overfladiske Iagttager ser ofte kun en ydre, mere eller mindre sprukken Skorpe. Selv hvor en saadan Skorpe har dannet sig, kan den gennembrydes i de energiske Villiesafgørelser, i hvilke det egentlige Jeg (le moi fondamental) ytrer sig, dette egentlige Jeg, der saa ofte forbliver skjult, men som betinger Sammenhængen i vor Personlighed. Kun den Handling bør kaldes fri, som har dette Jeg til Aarsag. I saadanne Handlinger virker hele Sjælelivets Forhistorie med; først og fremmest stiger da det inderste Selv (le moi d’en bas) \nm{39} op til Overfladen og gør sig gældende i Modsætning til Vaner og til Tilfældigheder. Saa pludselige og overraskende saadanne Afgørelser kunne være, vil dog nøjagtig Iagttagelse vise, at der er foregaaet en Gæringsproces helt nede i Dybet, --- at en tiltagende Spænding i Følelses- og Tankelivet er gaaet forud, hvorved de tilsyneladende motivløse Handlinger finde deres Forklaring. Vor Karakter er os selv, og vore Handlinger ere frie, naar de udspringe af den. Frie Handlinger i denne Betydning ere sjeldne, og Friheden har mange Grader. De fleste af vore daglige Handlinger skyldes øjeblikkelige Tilskyndelser eller fæstnede Vaner; den store Energiudfoldelse spares til de afgørende Situationer --- og selv hvor hele vor Personlighed burde vibrere, kommer ofte Inerti og Magelighed til at raade.

Bergson’s Opfattelse af Begrebet „Frihed“, anvendt paa Villien, falder ganske sammen med den Definition, Spinoza har givet: „Det Væsen kaldes frit, der existerer ud af sin egen Naturs Nødvendighed og bestemmes til at handle af sig selv alene.“ Frihed i denne Betydning danner ikke Modsætning til Nødvendighed.

Alligevel polemiserer Bergson baade mod „Determinismens Vildfarelse“ og mod „dens Modstanderes Illusion“. Deterministernes Vildfarelse skal bestaa i, at de gøre en udvortes,
% --- p. 46 ---
\opage{46}mekanisk Sondring mellem Aarsag og Virkning og lade de enkelte Motiver optræde overfor hverandre som isolerede og indbyrdes uafhængige Elementer. Det er den engelske „Associationspsykologi“, han her tænker paa; men hans Kritik af den rammer ikke enhver Opfattelse, som antager Aarsagslovens Gyldighed for Villieslivet. Denne Gyldighed antager Bergson selv, og det i selve sin Anvendelse af Begrebet Frihed som Udtryk for Karakterens indre Sammenhæng. Ja, han udtaler, at det er farligt at definere Begrebet „Frihed“, da man saa ikke undgaar at blive Determinist, og han ser ikke, at han faktisk selv har defineret det ved sin Anvendelse af det. --- Indeterministernes „Illusion“ finder Bergson deri, at de antage absolute Begyndelser af Villiesakter og derved løsrive disse fra Sammenhængen med hele Sjælelivet. Det er her stadigt den samme Fjende, han bekæmper, som overalt i hans Filosofi.

Bergson har naturligvis Ret i, at menneskelige Handlinger ikke kunne forudsiges med samme Sikkerhed som f. Ex. astronomiske Fænomener. Og han sætter, ligeledes med Rette, dette i Forbindelse med, at Gentagelse i langt mindre Grad kan paavises paa det aandelige end paa det materielle Omraade. Dertil kommer, at vi for en stor Del kun kende og forstaa de foregaaende sjælelige Tilstande gennem deres Indflydelse paa de følgende; især naar vi skulle opfatte Andres Sindstilstande ere vi væsentligt henviste til deres Resultater. Vi ere uheldigere stillede end Romanforfatteren, der forud kan vide, til hvilket Resultat hans Karakterers Udvikling fører; vi maa selv gennemleve de Andres Sindstilstande for at forstaa deres Betydning --- og hvorledes er dette muligt, da vi og de have hver sine særskilte Personligheder? Tilmed vilde der, selv om det kunde lykkes, ikke blive Tid til nogen Forudseen, da vi ikke maa afkorte, men maa gennemleve Alt, lige til Afgørelsens Øjeblik kommer.

26. Ogsaa hvad et helt andet psykologisk Fænomen angaar, hænger for Bergson Forstaaelsen sammen med hans filosofiske Hovedtanke om Modsætningen mellem det virkelige indre Liv og de ydre Former og Vaner, som praktiske og so%
% --- p. 47 ---
\opage{47}ciale Forhold medføre. Det er Latteren og det Komiske, som han har gjort til Genstand for et særeget lille Skrift (\textit{Le rire. Essai sur la signification du comique.} 1900).

Mellem Naturen og os, ja, mellem vort eget Sjæleliv og os selv, lægger der sig ligesom et Slør, fordi vi maa arbejde for Selvopholdelsen og fyldestgøre Samfundets almene Regler. Vi nødes til at lægge Vægt paa enkelte ydre Egenskaber og Forhold ved Tingene og ved os selv. Og Sproget forøger ogsaa her Udvortesheden og Mekanismen ved at danne Ord for hver Enkelthed. Vi komme ikke Ansigt til Ansigt med Tingenes og Væsenernes virkelige, individuelle Beskaffenhed. Vi se ikke Tingene selv, men lade os fordetmeste nøje med at læse de Etiketter, der ere fæstede paa dem. Og dette gælder ikke blot om de ydre Genstande; det gælder ogsaa vore egne Sindstilstande, som unddrage sig os, netop hvad det mest Inderlige, Personlige, oprindeligt Oplevede angaar.

Men heldigvis formaar Kunsten at hæve dette Slør og vise os Tilværelsen i dens hele Oprindelighed. Ti Kunsten er et Værk af Sjæle, som have formaaet at frigøre sig fra den praktiske Trang og dens Virkninger. Det er især den store Tragiker, som blotter Sjælelivets Inderste for os. Ved Hjælp af sine egne Oplevelser formaar han at læse i de andre Menneskers Indre. Han behøver ikke direkte at have oplevet de Stemninger og Handlinger, han skildrer. Men han har hos sig selv erfaret Muligheder, der, hvis de naaede frem til Virkelighed, vilde føre til saadanne Karakterer, som han skildrer. Derved blive disse Karakterer Dele af hans eget Væsen, og han skuer som i en Vision Virkeligheden i dens fulde og hele Individualitet.

Hvad der i den tragiske Kunst naas direkte og positivt, naas i den komiske Kunst indirekte og negativt.

Komikeren bygger væsentligt paa Iagttagelse af andre Mennesker. Hos sig selv vilde han ikke finde noget Latterligt, selv om han søgte derefter; ti vi ere kun latterlige ved den Side af vor Personlighed, som unddrager sig vor Bevidsthed. Og hos de Andre er det kun den ydre Adfærd, Overfladen, han kan bruge. Og den er mere eller mindre fælles for flere
% --- p. 48 ---
\opage{48}forskellige Individer. Hvis Komikeren trængte ind i sine Modellers inderste Sjæleliv, --- hvis han naaede ind til den indre Bevægelse, som foregaar under Skorpen, skønt stadigt hæmmet af denne, --- vilde det være forbi med Komiken. Komiken har sit Væsen og sin Betydning som en Reaktion mod alt det Vanemæssige, Mekaniske, Upersonlige i Menneskelivet. Den er derfor ikke ren Kunst, har ikke den egentlige Kunsts uinteresserede Karakter. Der er stadigt en Bagtanke, som gaar ud paa at rette og oplyse.

Det er en social Gerning, Komiken udøver. Den har stedse en social Forudsætning. Latteren er altid en Gruppes Latter. Det Latterlige beror paa Modsætningen til et større eller mindre Samfunds Sæder og Forestillinger. Nu er det ganske vist for en stor Del Samfundet selv, der frembringer den Stivnen, den Skorpe, der bliver Latterens Genstand. Samfundet kræver jo bestemte ydre Former og Regler overholdte af Alle og kan derved komme til at gøre Mennesker upersonlige. Men paa den anden Side kan Samfundet ikke staa sig ved, at Menneskene blive blotte Maskiner og Vanevæsener, saa at Selvvirksomhed falder bort, og nye Opgaver, Arbejde i nye Retninger hindres. Derfor virker Komiken i Samfundets Tjeneste og er et socialt Redskab. ---

Man ser, hvorledes Bergson’s Hovedtanke kommer igen i hans Teori om Latteren som den, der blotter Misforholdet mellem Mekanisme og Personlighed, mellem Skal og Kærne, mellem det Stive og det Bølgende.

27. Der er endnu et Moment i det Latterlige, som Bergson nævner, men som han efter min Opfattelse ikke lægger tilstrækkelig Vægt paa, skønt det er et Hovedpunkt. Det er, at det Latterlige forudsætter en pludselig Afbrydelse af en Sammenhæng. Latteren som Kontrastvirkning kommer ikke til sin Ret. Kontrasten er efter Bergson kun en Betingelse, der tjener til at vække Opmærksomheden. Erfaringen viser dog, at der maa gaa en mere eller mindre bevidst Forventning forud, som pludseligt opløses i Intet. Altsaa Kontrasten virker i
% --- p. 49 ---
\opage{49}nedadgaaende Retning --- i Modsætning til det Ophøjede, hvor den pludselige Kontrast ogsaa virker, som hvor Overgangen sker i opadgaaende Retning, fra det Smaa til det Store.

Bergson’s egen Teori omhandler kun et enkelt Exempel paa saadan nedadgaaende Kontrast. Vi vente, naar et Menneske træder os imøde, at træffe noget Levende, noget Individuelt, noget Nyt; og naar vi saa istedetfor træffe en Maskine, en død Gentagelse af Fortiden, den egne eller Andres, kan denne Skuffelse udløse Latteren.

Ogsaa nogle af Bergson’s mere specielle Exempler pege i samme Retning. Naar Ceremonier uden Indhold eller naar Gentagelser virke latterligt, saa er det, fordi en Forventning opløses i Intet. Det bedste Exempel er maaske her den Maade, paa hvilken Lighed mellem Personer kan virke latterligt, Noget, som Pascal allerede havde gjort opmærksom paa, og som Bergson med Rette tager til Indtægt for sin Teori. „Det virkeligt levende Liv,“ siger han, „vilde aldrig gentage sig. Hvor der er Gentagelse, fuldstændig Lighed, formode vi stedse en Mekanik virkende bag det Levende. Analyser det Indtryk, man faar overfor to Ansigter, der ligne hinanden for meget; man vil finde, at man kommer til at tænke paa to Exemplarer, der ere støbte i samme Form, eller to Aftryk af samme Signet, eller to Reproduktioner af samme Cliché, idethele paa en industriel Fabrikation. Denne Livets Bøjning i Retning af Mekaniken er her den sande Grund til Latteren.“ Men dette Exempel peger ud over Bergson’s Teori. Ti det er i saadanne Tilfælde ikke denne Ideassociation, der er det Afgørende. Det Afgørende er selve Overgangen fra det ene Ansigt til det andet; det er Forventningen om et nyt Menneske, som skuffes. Pascal lægger med Rette Vægt paa, at hvert Ansigt for sig ikke behøver at vække Latter; det er selve Sammenstillingen, der bliver Aarsagen.

28. Ogsaa i andre Henseender viser det sig, at Bergson’s Teori kun passer paa en enkelt Gruppe af Tilfælde.

Han kender kun Latter som ydmygende, som Udtryk for
% --- p. 50 ---
\opage{50}Bitterhed og Haan. „Den, der ler,“ siger han, „hævder sig selv paa mere eller mindre stolt Vis og tenderer til at betragte den Anden som en Marionet, hvis Traade han selv holder. I denne Selvoptagethed ville vi hurtigt opdage en Smule Egoisme, og bag Egoismen Noget, som er mindre uvilkaarligt og mere besk.“ --- Men vi le jo dog ogsaa ved pludselig Befrielse og Lettelse, selv om det ikke gaar ud over nogen Anden, naar blot en Hindring eller et Tryk viser sig i sin Intethed. Og vi kunne le i Sympati med selve den, vi le ad, maaske endog i Sympati med selve det Træk eller den Handling, der fremkalder Latteren. Vi kunne ogsaa le ad os selv, og i saadanne Tilfælde viser det sig da, at Bergson ikke har ganske Ret i den Sætning, at vi kun ere latterlige ved det i vort Væsen, der undgaar vor egen Opmærksomhed. Bergson omtaler kun Haanens, ikke Humorens Latter.

Og naar han hævder, at Latteren stedse er en Gruppes Latter, saa maa dertil bemærkes, at det i hvert Tilfælde ikke behøver at være en faktisk given Gruppe, der forudsættes. Latteren kan forudsætte en Appel til Sandhed, Ærlighed og Retfærdighed, selv om disse Ideer kun kunne tænkes realiserede i et idealt Samfund, og den kan da være Udtryk, snart mere for en haanende, snart mere for en humoristisk Stemning hos En, der staar ene i sin Tid og dog ikke føler sig ene \nm{40}.

Det er derfor en urigtig Modsætning, i hvilken den komiske Kunst i Bergson’s Teori kommer til at staa til den tragiske Kunst. Giver hin os ikke, ligesaa vel som denne, Livets Realitet? Livet er, som allerede Platon saa, baade at kalde en Komedie og en Tragedie, begge Dele paa én Gang. Og naar Kunsten skal afsløre Virkeligheden, der stedse er individuel, saa maa man tage Virkeligheden baade i dens Begrænsning og i dens Storhed, baade i dens Udvorteshed og i dens Inderlighed. Al individuel Existens omfatter begge Dele, og derfor udfordres der en dobbelt Kunst til at hæve det Slør, der kan dække den for os. ---

Naar iøvrigt Bergson lærer, at „Kunsten ikke har andet Maal end at fjerne de praktisk nyttige Symboler, de konven%
% --- p. 51 ---
\opage{51}tionelt og socialt accepterede Generaliteter, med et Ord, Alt hvad der maskerer Realiteten for os, for at bringe os Ansigt til Ansigt med selve Realiteten,“ --- saa fører den Betydning, han herved tillægger Kunsten, naturligt til at mindes den Betydning, han tillægger Filosofien i Modsætning til den analyserende og abstraherende Videnskab, og vi komme til at spørge om Forholdet mellem Kunst og Filosofi. Derved føres vi endnu engang tilbage til Begrebet Intuition.

\chapterend

% --- p. 52 ---
\chapterhead{FORSØG PAA METAFYSIK}

\opage{52}29. \initial{A}{rbejdsdelingen} trænger sig frem overalt, ogsaa i Filosofien. De forskellige filosofiske Problemer fjerne sig mere og mere fra hverandre, og der vil maaske tilsidst kun kunne paavises en fælles Type for dem, en Type, hvis Udspring saa stiller et nyt Problem. Psykologi, Erkendelsesteori og Etik tage hver for sig deres særegne Udgangspunkter, og det kommer da an paa, om Tankerækkerne ud fra disse Udgangspunkter kunne mødes i Harmoni. For Bergson’s Vedkommende er det intet af de for de nævnte Fag ejendommelige Problemer, der staar som det centrale, men derimod et Problem, der er traadt i Baggrunden for Filosoferne af den kritiske og den positivistiske Retning, ikke fordi det nogensinde skulde kunne falde bort, men fordi det er forbundet med særegen Vanskelighed, og fordi dets Behandling forudsætter en Behandling af de nævnte Fag. Det er Tilværelsesproblemet (det kosmologiske eller metafysiske Problem). Bergson giver sig dristigt i Kast med det. Ligesom Faust er han utilfreds med, hvad der ellers kan vides, naar han ikke kan løse Verdensgaaden. Han ønsker som Faust, at han „den Kraft kan fatte og finde, der virker i Verden inderst inde, den sidste Løsning paa alle Gaader“. For at naa dette Maal hengav Faust sig til Magien; Bergson paakalder Intuitionen. Det er ogsaa ved dens Hjælp, han søger at løse den Modsætning, hans Hovedtanke om Modsætningen mellem Liv, særligt personligt Liv, og Mekanisme stiller. I Intuitionen ere vi ude over denne Modsætning.
% --- p. 53 ---
\opage{53}I et Foredrag paa Filosofkongressen i Bologna 1911 udtalte Bergson: „En Filosof, som fortjener dette Navn, har aldrig sagt mere end en eneste Ting, og denne ene Ting har han snarere søgt at sige, end han virkeligt har faaet den sagt. Og han har kun sagt en eneste Ting, fordi han kun har set et eneste Punkt, og saa var det endda snarere en Berøring end et Syn!“ \nm{41} Det er sikkert Bergson’s egen Erfaring som Tænker, der har lagt ham disse Ord i Munden. Alt hvad han har skrevet, drejer sig, som vi have set, om én Hovedtanke og om Løsningen af det Problem, den indeholder. Og man mærker gennem hans mange Vendinger og Billeder, i hans stadige Kredsen om det ene Punkt, og i Kampen med Sproget og med Begreberne, hvor det koster ham Anstrengelse at fastholde og udtrykke, hvad der for ham er det Centrale i Tankens Verden. Den værdifuldeste Side ved den store Indflydelse, han øver paa store Kredse, hænger sammen med denne Ejendommelighed hos ham. Forhaabentligt vil denne Side af hans Filosofi virke stærkere end den af Misforstaaelse udsprungne Konsekvens, man har draget af den, at man nu kunde hengive sig til drømmende Skuen, fordi videnskabeligt Tankearbejde ikke fører op til Højdepunktet. Bergson selv har, for at kunne fremstille sin Hovedtanke og sin Løsning fra alle Sider, maattet arbejde ihærdigt og strengt; for at naa sit Maal har han behandlet en Række af speciellere Spørgsmaal og kastet Lys over dem. Især har han som beskrivende Psykolog gjort en god Gerning. Men Psykologien er for ham kun det Grundlag, fra hvilket han stiger op til den metafysiske Intuition.

Nu er Intuitionen en Skuen, eller, som Bergson siger, et Syn, en Berøring. Og det er foreløbigt kun et Billede af vort eget inderste Liv, den bringer os. Hvorledes kunne vi da heraf udlede en Metafysik, som jo skal være en Lære, en Teori? Det er jo dog kun en Side eller en Del af Tilværelsen, Intuitionen kan vise os, og hvorledes naa vi derfra til en Opfattelse, som kunde gælde for Tilværelsen idethele?

Dette Spørgsmaal stiller Bergson sig ikke klart og bestemt, og det er dog selve Metafysikens Mulighed, det gælder. Saa Meget er dog klart, at han ikke, saaledes som Hegel, gaar
% --- p. 54 ---
\opage{54}Konstruktionens, den dialektiske Udlednings Vej. Men hvilken anden Vej er der da, der kan føre fra Delen til Helheden? Der gives neppe nogen anden Vej end den, allerede Leibniz henviste til, naar han hævdede, at al Metafysik grundede sig paa Analogi. Til samme Resultat kom ogsaa Kant gennem en mere kritisk Fremgangsmaade. Og det er i Virkeligheden, som jeg allerede paaviste i min Bog „Den menneskelige Tanke“, den Vej, ogsaa Bergson gaar.

Med Urette bebrejder Bergson de filosofiske Systemer, at de uden videre erklære, at en Del af Tilværelsen repræsenterer den hele Tilværelse. Det maa ske saa i ethvert System, enhver Metafysik, og det sker ogsaa i Bergson’s. Hans Intuition viser os vort indre Liv som en stadig Udfoldelse, en Stræben, en Nyskaben. Hvad vi saaledes finde som det Inderste i den Virkelighed, vi selv ere, bliver en Type, et Mønster (modèle), efter hvilket vi forestille os al anden Virkelighed, naar det lykkes os at hæve os over den praktisk nødvendige, men udvortes Maade at se Tingene paa. Al anden sand Virkelighed kan da kun være gradsforskellig fra den, Intuitionen viser os i os selv. Det er overalt en Stræben, en Higen, et Opsving (élan), der ligger til Grund. Tilværelsens Kærne er ikke død Væren, men Vorden, Udfoldelse, Villen (tendance). I „Évolution créatrice“ har Bergson, efter sit eget Udsagn, villet vise, „at det Hele er af samme Natur som Jeget, og at man griber det ved en stedse mere fuldstændig Fordybelse i sig selv.“ \nm{42}

Bergson’s Metode er da ikke blot intuitiv; Intuitionen er kun det første Skridt, Analogislutningen maa gøre de følgende Skridt. Analogier kunne være højst frugtbare for Forskningen; de kunne føre til Opdagelser og til nye Synsmaader; men bevisende ere de ikke \nm{43}, og navnligt ikke, naar de ikke skulle tjene til at belyse et Iagttagelsesomraade ved Hjælp af et andet, men naar de skulle forklare Tilværelsen overhovedet, hvor ingen Erfaringsbekræftelse, ingen Efterprøvelse af Analogiens Gyldighed er mulig. Alle religiøse og metafysiske Forestillinger bygge tilsidst paa mere eller mindre dristige Analogier; men saadanne Forestillinger overskride da ogsaa Grænsen mellem Tankens og Poesiens Verden.
% --- p. 55 ---
\opage{55}Lad os se noget nærmere paa, hvorledes Bergson anvender sine Intuitioner og Analogier.

30. I sin Naturopfattelse stiller han, som vi have set, Livet i en skarp Modsætning til det Livløse. Han betoner Livets historiske Karakter. Det Nye hænger kontinuerligt sammen med det Gamle; men det kan derfor meget vel træde i uventet Modsætning til det, saa at Forudsigen bliver en Umulighed. De fysiske og kemiske Love indeholde kun Betingelser og Skranker for Livet, angive ikke selve Livets Aarsag. Og da den exakte Naturvidenskab mere og mere anvender en rent matematisk Behandlingsmaade, medens paa den anden Side Livet, naar vi ret trænge ind i det, frembyder en Analogi, der maaske er mere end en Analogi, med vort Sjæleliv, --- saa er det tilsidst Matematik og Psykologi, der komme til at staa overfor hinanden i Striden om Verdensanskuelsen. Matematisk Naturvidenskab hviler paa den Forudsætning, at de samme Elementer og de samme Forhold atter og atter gentage sig; Psykologien viser os derimod en Væxt, en Nydannelse, der med Rette kan kaldes en Nyskaben. Begreberne Væxt og Skabelse ere for Bergson Et: l’idée de création se confond avec celle d’accroissement. Direkte kende vi saadan Væxt eller Nyskaben fra vort eget Sjæleliv, men den er alt Levendes Særkende under dets stadige Kamp mod Materiens Tendens til at ende i absolut Ligevægt.

„Naar jeg“, siger Bergson \nm{44}, „betragter den Verden, i hvilken vi leve, finder jeg, at denne sammenhængende Helheds automatiske og strengt bestemte Udvikling er en Virken, som aftager (de l’action qui se défait), og at de uforudsete Former, som Livet afgrænser indenfor det, Former, der selv kunne forlænge deres Existens gennem uforudsete Bevægelser, betyde en Virken, der er i Væxt (de l’action qui se fait). Nu har jeg al Grund til at tro, at de andre Verdener ere analoge med vor, --- at Tingene foregaa der paa samme Vis. Jeg véd, at de ikke alle ere blevne til paa samme Tid, da Iagttagelse endnu den Dag idag viser mig Stjernetaager ifærd med at koncentrere sig. Dersom det overalt er den samme Art af Virken, der gaar for sig, hvad enten den aftager eller voxer, saa kan jeg ganske
% --- p. 56 ---
\opage{56}simpelt udtrykke denne sandsynlige Lighed ved at tale om et Centrum, hvorfra Verdenerne vælde frem, ligesom Raketter ved et uhyre Fyrværkeri, --- kun at jeg ikke opfatter dette Centrum som en Ting, men som en stadig Fremvælden (une continuité de jaillissement). Naar man definerer Gud paa denne Maade, er der intet en Gang for alle Færdigt i hans Væsen; han er uophørligt Liv, Virken, Frihed. Opfatter man Skabelsen paa denne Maade, er den intet Mysterium; vi erfare den i os selv, naar vi handle med Frihed.“ „Det gælder blot om at skue for at skue, og ikke mere skue for at handle. Da aabenbarer det Absolute sig meget nær ved os, og paa en Maade i os. Det er af psykologisk, ikke af matematisk eller logisk Natur. Det lever med os. Det lever i en psykologisk Tid (il dure), ligesom vi gøre, kun i flere Henseender mere koncentreret, mere samlet i sig selv.“

Ud fra, hvad der er grebet i Intuitionen, stiger, som man ser, Bergson ned ad Analogiens Skala til den materielle Natur, hvis Virken er en Tendens til Ligevægt, og op ad til et Verdenscentrum som Udtryk for det mest koncentrerede Liv. Fysik, Biologi, Psykologi og Teologi samle sig i denne omfattende Intuition i et Grundsyn. Dette er Højdepunktet i Bergson’s Filosofi.

31. Jeg gaar nu tilbage for at undersøge det Grundlag, ud fra hvilket Bergson er naat til dette Højdepunkt. Det væsentlige Punkt er her udtrykt i den Anklage, han retter mod al Videnskab og al Filosofi, den moderne saavel som den antike, at de tro at kunne forklare Bevægelse ud fra Hvile, Forandring ud fra det Uforanderlige, det Levende ud fra det Livløse. Der har her efter hans Mening udviklet sig en intellektuel Vane, fra hvilken vi kun kunne frigøre os ved en Villiesakt, ved et Spring over i en umiddelbar Skuen, der lader os mærke Livets Puls, som den banker ikke blot i os, men i den hele Tilværelse. Bergson mener, at Grækeren Zenon’s Opfattelse af Bevægelsen som resulterende af ubevægede Tilstande stadigt er klassisk og typisk for Videnskaben, idet denne stedse maa opfatte Tiden som Sum af enkelte udelelige Øjeblikke, ligesom den opfatter Naturen som en Sum af udelelige Atomer.
% --- p. 57 ---
\opage{57}Heri har han dog afgjort Uret. Den antike Opfattelse var ganske vist tilbøjelig til at sætte det Uforanderlige over det Foranderlige og til at mene, at der om det Foranderlige ingen sand Videnskab kunde være. Men den nyere Tænkning og Videnskab tog, med Bruno og Galilei, afgjort sit Stade i Forandringens Verden og stillede sig til Opgave at finde Love for Bevægelse og Udvikling. Selv Hvilen blev opfattet som potentiel eller minimal Bevægelse. Alle Naturlove, den nyere Videnskab har paavist, ere Love for Bevægelse, og den store Rolle, Udviklingsbegrebet spiller, peger ogsaa i helt anden Retning end den antike Uforanderlighedslære. De Elementer eller faste Former, man foreløbigt stanser ved, betegne kun Stationer, Hvilesteder, ikke uforanderlige Stader. Hverken Mangfoldighed eller Enhed, hverken Diskontinuitet eller Kontinuitet kunne udtrykke Tankens Virken og Tilværelsens Natur; de ere hver for sig kun Abstraktioner. Det virkelige Tankearbejde, saavel som Alt, hvad vi kunne formode om Tilværelsens Natur, frembyder en stadig Brydning mellem Kontinuitet og Diskontinuitet og stadige Svingninger imellem dem, uden at det ene af disse Synspunkter kan udledes af det andet.

32. Grunden til den Vægt, Forstanden lægger paa uforanderlige Elementer, finder Bergson i, at den staar i det praktiske Livs Tjeneste. I sin Opfattelse af den menneskelige Erkendelse er han afgjort Pragmatiker. \nm{45} Men han bryder med Pragmatismen, og med den videnskabelige Betragtningsmaade idethele, naar han gaar over til sin egentlige Filosofi, der kræver Intuition og sætter umiddelbar Anskuen som det Højeste. Det er for ham Filosofiens Opgave at føre ud over det Diskontinuerlige og Splittede til en Skuen.

Men hvorledes bliver dette muligt, naar man, som Bergson gentagne Gange udtaler \nm{46}, vel kan gaa fra Intuition til Forstand og Analyse, men ikke omvendt, ligesom man kun skal kunne gaa fra Bevægelse til Hvile, men ikke omvendt? --- Man maa kunne gaa begge Veje, gøre Intuitionen tydeligere gennem Analyse og samle det, Analysen sondrer, gennem ny Intuition. Som vi ovenfor have set, staa Intuitioner af første, tredje og fjerde
% --- p. 58 ---
\opage{58}Art (a, c og d) i stadig Vexelvirkning med hverandre. Bergson ender i en Dualisme mellem Intuition og Analyse, som kun skal kunne overvindes ved den koncentrerende Villiesakt, han appellerer til, og som han ved sin Fremstillingskunst søger at fremlokke.

Naar det lykkes os at naa til en saadan Intuition, faar ifølge Bergson Alt Liv og Bevægelse for os; alt Dødt og Uforanderligt forsvinder; vi føle os dragne, løftede, baarne. „Vi \textit{leve} mere end før, og dette Overskud af Liv fører den Overbevisning med sig, at de alvorligste filosofiske Gaader kunne løses, eller maaske endog, at de ikke burde opstilles, at de udspringe af et stivnet Blik paa Universet.“ \nm{47}

Men der bliver da et nyt Problem til, saasnart vi vaagne op af Intuitionen: hvorledes forholder den problemløse Tilstand sig til den af Problemer betagne? Hvorfor er der Forskel mellem Anskuen og Forstand, mellem Overbevisning og Begrundelse? I selve Intuitionens Tilstande glemmer man, at der existerer aandelige Tilstande af helt anden Art; men disse ville atter og atter gøre sig gældende, og de finde ikke deres Forklaring ved selve Intuitionen, som jo ser sin Fjende i dem!

Forskellen mellem de to Tilstande, den, i hvilken Problemer rejse sig, og den, i hvilken de give Plads for umiddelbar eller genvunden Anskuen, peger tilbage til en Enhed, naar man agter paa, at den menneskelige Aands Natur lægger sig for Dagen i begge Tilstande. Bergson og Pragmatismen have Uret, naar de ikke kunne finde Personlighedens Væsen i videnskabelig Forsken ligesaa vel som i den umiddelbare Skuen. Selve de Metoder, Videnskaben udformer, og de praktiske Veje, Livet nøder os til at gaa, betinges ved den menneskelige Aands Natur og vilde være uforstaaelige uden den. Allerede rent biologisk lægger dette sig for Dagen i det nøje Forhold, der er imellem Organ og Funktion. Erkendelsesteoretisk og psykologisk ere vi nu saaledes stillede, at vi kun kende Organet gennem Funktionen. Det er gennem Studiet af den Maade, hvorpaa menneskelig Forsken og Handlen fremtræder, at vi kunne lære den menneskelige Aands Natur at kende. Vi afskære os fra saadant
% --- p. 59 ---
\opage{59}Kendskab, naar vi paa Forhaand stille os udenfor Handlingens og Forskningens Liv, som jo dog ogsaa hører med til Livet \nm{48}.

Og paa hvilket Sprog skulde Intuitionens Indhold kunne udtrykkes, naar Intuitionen staar i absolut Modsætning til Forstanden, og „kun Forstanden har et Sprog“? \nm{49} Forstanden kan jo ikke hjælpe os, da den er henfalden til Trældom i det praktiske Livs Tjeneste. Som vi have set (3), er Bergson klar over dette Punkt og søger gennem Sammenligninger, der hver for sig halte, men alle pege i samme Retning, at fremkalde Intuitionen indirekte, at fremlokke den Villiesanstrengelse, den Samling, som Opgaaen i Skuen forudsætter. Og her staar en glimrende Fremstillingskunst til hans Raadighed. Men ikke blot det. Han har Ret i, at enhver dybere gaaende aandelig Paavirkning, paa det intellektuelle Omraade som paa alle andre, egentligt talt altid sker indirekte, gennem Vækkelse til Selvvirksomhed. Og ogsaa i Videnskaben kan det være af Betydning, at et Totalbillede, en samlet Skuen af en Sammenhæng bliver til. Det er Fantasiens Værk i Videnskabens Tjeneste. Men den videnskabelige Fantasi adskiller sig netop derved fra den kunstneriske, at den ikke er en absolut Afslutning, men stedse tjener som Grundlag og Indledning til ny Eftertanke, medens det kunstneriske Billede har sin Værdi, sit Formaal i sig selv. Derfor bliver Bergson’s Intuition Kunst, ikke Videnskab eller Filosofi. Han anfører gentagne Gange \nm{50} Kunsten som Analogi til den filosofiske Intuition. Men det bliver mere end Analogi, naar Intuitionen træder i den skarpe Modsætning til Eftertanken, som Bergson antager. Og om en Løsning af Problemer er der i hvert Tilfælde ikke Tale; derimod sker der, som alt bemærket, ved Overgangen til Intuition en Overgang til en problemløs Tilstand.

Det er tilsidst Modsætningen mellem Kunst og Videnskab, der betegner den Bergsonske Filosofis Karakter. Det er som en Modsætning mellem Hellig og Søgn, mellem Poesi og Prosa. Men saaledes er nu engang Livet: det bevæger sig gennem Modsætninger, og en tapper Tanke ser disse Modsætninger under Øjne, netop fordi de begge høre med til Livet. Det gælder om i Prosaen at finde Poesiens skjulte Aare, og i Poesien ikke
% --- p. 60 ---
\opage{60}at glemme, at Vejen igen vil føre tilbage til Livets Prosa. Der gives en Tankearbejdets Poesi, som Pragmatismen er tilbøjelig til at nægte, og som oftest virker med en indre besjælende Magt, der kun paa Højdepunkterne frigøres og søger at finde Udtryk. Men den oplader sig kun for den trofaste Arbejder.

33. Paa de Punkter, hvor Intuitionen ret udfolder sig indenfor Bergsons Filosofi, overskrides ikke blot Grænsen mellem Tænkning og Poesi, men ogsaa Grænsen mellem Filosofi og Teologi. Som vi have set (30), er for Bergson Gudsbegrebet Et med den Kilde eller det Centrum, hvorfra Verdensudviklingen uafbrudt bryder frem. Paa alle Omraader, i alle Verdenssystemer, skønt mest direkte i Livsudviklingen, finder han et saadant Udspring for Virken og Leven. Den Modsætning, der ellers spiller saa stor en Rolle i hans Filosofi, Modsætningen mellem Liv og Materie, træder i saadanne Intuitionsøjeblikke tilbage for ham. Den hører da til de Problemer, der for den højeste Intuition vise sig uberettigede og ubegrundede (32).

En Filosof bør altid være forsigtig med at benytte teologiske Udtryk. Til de teologiske Forestillinger svare indenfor Filosofien i det Højeste Problemer, og man behøver kun at tænke lidt over Gudsbegrebet (hvad man ikke gør i Intuitionens Øjeblikke) for at se, at det indeholder det største Problem af alle, Problemet om, hvorvidt det sidste Grundlag for Værdi, for det Gode, kan tænkes absolut identisk med det sidste Grundlag for den Virkelighed, Erfaringen viser os. Dette store Problem har man ikke Ret til at skyde tilside. Og det skærpes ganske særligt derved, at Værdiernes Række (Rækken af Midler og Formaal) og Virkelighedens Sammenhæng (Rækken af Aarsager og Virkninger) hver for sig ere uafsluttelige. Der kan stedse tænkes højere Værdier, længere tilbage liggende Aarsager og længere ude i Fremtiden liggende Virkninger. Overfor de Problemer, der her stadigt stille sig paany, kan enhver Intuition kun betyde et foreløbigt Hvilested. Det kan kun være en Tro, --- en Tro, hvis psykologiske Natur og Betingelser Filosofien har at undersøge, --- naar der antages en Harmoni mellem Værdiernes og Aarsagernes Rækker \nm{51}.
% --- p. 61 ---
\opage{61}Bergson har paa Intuitionens Højdepunkt ment at finde en Enhed af Fysik, Biologi, Psykologi og Teologi (30). Han mener at have Ret til at bruge Ordet Skabelse om den stadige Frembringen af nye Fænomener paa de forskellige Omraader. Men Skabelse betyder for ham det Samme som Væxt, og han tænker sig Verdensudviklingen i Analogi med Strømmens Udspring af Kilden. Og en Kilde er jo selv en Strøm, --- den første Del af den Strøm, der siges at udspringe af den!

I sit Svar paa en Kritik, der fremkom i et af Jesuiter redigeret teologisk Tidsskrift, hævdede Bergson alligevel, at hans Filosofi begrunder „Tanken om en fri og skabende Gud“ (l’idée d’un dieu créateur et libre). Hertil blev der med fuld Ret svaret, at for den katolske Teologi er Gud ikke blot Kilden, af hvilken Strømmen vælder ud; Gud udfolder sig ikke til en Verden, men han frembringer den ved en Skabelse i helt anden Betydning af Ordet end den, i hvilken Bergson tager det \nm{52}.

Filosofien staar sig bedst og gør Aandslivet den bedste Tjeneste ved at gaa sine egne Veje og behandle ogsaa det religiøse Problem efter sine egne Metoder. Saa kan Regnskabet klart gøres op, og det kan ses, hvor langt man paa Grænseomraaderne kan komme ad rent menneskelig Vej. Arbejdet i denne Retning maa stedse tages op paany. Afslutning paa Tankearbejdet kan ikke være berettiget, saalænge Erfaringen stedse kan bringe nye Fænomener, Situationer og Problemer. En „Afrunding“ vil derimod naturligt ske fra Tid til anden, et foreløbigt Hvilested til Orientering, før der arbejdes videre. Men en saadan Afrunding er ikke Et med Afslutning. Der vil stedse være Noget udenfor den Kreds, indenfor hvilken vi runde af. Her gælde Goethes Ord: „Nie geschlossen, oft „geründet“.

Bergson’s „Intuition“ adskiller sig derved fra tidligere Forsøg i samme Retning, at den er en Skuen af Bevægelse, Forandring, Liv, Udvikling. Og han betoner selv stadigt, at ingen Enkelt kan føre os til de Intuitioner, i hvilke Tanken kan naa sin Afslutning; der udkræves Arbejde af mange Tænkere og fra mange Sider. Og desuden ere jo for ham de værdifuldeste Intuitioner de, der ere gaaede gennem Forstan%
% --- p. 62 ---
\opage{62}dens Skærsild, skønt der derefter maa gøres et Spring for at komme over i Intuitionen.

Enhver Filosof af Betydning er baade et Symptom paa Retninger i Tiden og en Faktor, der bestemmer saadanne Retninger. Paa Grund af det kunstneriske Element i Bergson’s Filosofi faar den maaske størst Interesse ved at være et bevidst Udtryk for en Strømning i vor Tid, der er udsprungen af Utilfredsstillethed ved, hvad rationel Tænkning og Erfaring hidtil have bragt. Ved sin glimrende Fremstillingskunst og sin utrættelige Grublen har han bragt denne Tendens frem for Dagen i principiel Form. Selv den, der mener, at Tankens Magt og Ret ikke har faaet den Plads og den Anerkendelse, der tilkommer den, kan fuldt ud beundre den Energi, hvormed den evige Grænsestrid mellem Livet og Videnskaben her er stillet Tidsalderen for Øje. Hvorledes end Udfaldet bliver af den Drøftelse, Bergson har givet Stødet til, og af den Kamp mellem Beundring og Forkastelse, hans Skrifter have fremkaldt, vil der sikkert være vakt en ny Interesse for Eftertanke over de store Spørgsmaal, Filosofien til alle Tider sysler med.

Jeg har naturligvis maattet opfatte og bedømme Bergson’s Filosofi ud fra mit eget Stade, men jeg har søgt at give en objektiv Karakteristik og at lade Kritiken paa naturlig Maade udvikle sig af den.

\chapterend

\clearpage
% --- p. 63 ---
\opage{63}\noterhead{NOTER}
\clearpage
\begin{small}

% --- p. 65 ---
\opage{65}\notelabel{1.} (p. 3.) Om Bergson’s Opfattelse af græsk Filosofi og dens stadige Eftervirkninger, se \textit{La Perception du Changement} (Oxford 1911) p. 6, og \textit{Évolution créatrice} (Paris 1907) p. 333--347. --- Om Parmenides se mine \textit{Mindre Arbejder.} Anden Række. p. 141--149, om Platon \textit{Mindre Arbejder.} Tredie Række. p. 82--101.

\notelabel{2.} (p. 4.) I Efteraaret 1905 gennemgik jeg paa Laboratoriet med en lille Kreds af Studerende „Les Données immédiates“. I Aaret 1907 gav jeg en kort Karakteristik og Kritik af Bergson i „Lehrbuch der Geschichte der neueren Philosophie“ (senere oversat paa Engelsk). I „Den menneskelige Tanke“ (1910), af hvilken der 1911 kom en fransk Udgave, drøftede jeg dels hans Opfattelse af det i psykologisk Forstand umiddelbart Givne (p. 8 ff.; 16 f.), dels hans Opfattelse af Villien (p. 292 f.), dels hans Stilling til Spørgsmaalet om Muligheden af en videnskabelig Metafysik (p. 315--318).

\notelabel{3.} (p. 5.) Se \textit{Den menneskelige Tanke} p. 56--58 om Forholdet mellem Anskuen og Dømmen.

\notelabel{4.} (p. 6.) \textit{La Perception du Changement} p. 5--6.

\notelabel{5.} (p. 7.) \textit{Introduction à la Métaphysique.} p. 4--9. (Revue de Métaphysique et de Morale. 1903).

\notelabel{6.} (p. 8.) Se Diskussionen i det franske filosofiske Selskab d. 28. Novbr. 1907. (Bulletin de la Société Française de Philosophie. VIII p. 18; 21).

\notelabel{7.} (p. 9.) De fleste af de i det Foregaaende nævnte Forfattere omtales i „Den nyere Filosofis Historie“ (se Registret) eller i „Moderne Filosofer“. \emph{Félix Ravaisson} er bedst kendt for sin Bog \textit{La Philosophie en France au XIX\textsuperscript{e} siècle} (Paris 1868), \emph{Lachelier} fra sit Skrift \textit{Du fonde%
% --- p. 66 ---
\opage{66}ment de l’induction} (1871) og sin Afhandling \textit{Psychologie et Métaphysique} (1885). Begge sidstnævnte Arbejder af Lachelier ere udkomne paany i et Bind af „Bibliothèque de Philosophie Contemporaine“ (1896). --- Af den i det Følgende omtalte \emph{Alfr. Fouillée} er i denne Sammenhæng især det Værk af Interesse, han udgav kort før sin Død: \textit{La Pensée et les nouvelles écoles anti-intellectuelles} (1910), som for en stor Del er rettet imod Bergson.

\notelabel{8.} (p. 12.) Smlgn. \textit{Moderne Filosofer}, p. 42--46. --- Smlgn. desuden med Hensyn til det hele Spørgsmaal \emph{Dwelshauvers}: \textit{La synthèse mentale.} Paris 1908.

\notelabel{9.} (p. 14.) Revue scientifique 1875. Tannery har først senere navngivet sig som Forfatter til denne Artikel. Smlgn. \emph{Bergson}: \textit{Les données immédiates.} p. 50 og Th. \emph{Ribot}: \textit{La Psychologie Allemande contemporaine} (Paris 1879). p. 200; 204 ff. --- Det følgende Citat i Texten af Bergson findes i nævnte Skrift p. 42.

\notelabel{10.} (p. 15.) \textit{Les données immédiates.} p. 98; 103.

\notelabel{11.} (p. 16.) Se om dette Punkt \textit{Den menneskelige Tanke}, p. 189--193. I sin Artikel \textit{Some Antecedents of the Philosophy of Bergson} (Mind. October 1913) har \emph{Lovejoy} paavist, at Bergson i sin Overseen af Forskellen mellem Begrebet Tal og Begrebet Kvantitet har en Forgænger i Ravaisson. Idet Lovejoy antager, at Bergson idethele er paavirket af Ravaisson, mener han gennem denne at spore en Indflydelse af den Schellingske Filosofi. Ravaisson havde været Schellings Tilhører i München.

\notelabel{12.} (p. 18.) \textit{Évolution créatrice}, p. 150 ff.; 179--182. --- Jeg anfører en karakteristisk Sætning: Si la conscience qui sommeille dans l’instinct se réveillait, s’il s’intériorisait en connaissance au lieu de s’extérioriser en action, si nous savions l’interroger et s’il pouvait répondre, il nous livrerait les secrets les plus intimes de la vie. (p. 179).

\notelabel{13.} (p. 18.) \textit{Évolution créatrice,} p. 192: C’est à l’intérieur même de la vie que nous conduirait \textit{l’intuition, je veux dire l’instinct devenu désintéressé, conscient de lui-même,} capable de réfléchir sur son object % PRINTED AS IS: object (objet)
et de l’élargir indéfiniment.

\notelabel{14.} (p. 20.) Et Hovedsted er \textit{Évolution créatrice} p. 198: La conscience se déterminant en intelligence, c’est-â-dire se concentrant d’abord sur la matière, semble ainsi s’extérioriser par rapport à elle-même; mais, justement parce qu’elle s’adapte aux objets du dehors, elle arrive à circuler au milieu d’eux, à tourner les barrières qu’ils lui opposent, à élargir
% --- p. 67 ---
\opage{67}indéfiniment son domaine. Une fois libérée, elle peut d’ailleurs se replier à l’intérieur, et réveiller les virtualités d’intuition qui sommeillent encore en elle.% PRINTED AS IS: "c’est-â-dire" (â for à)

\notelabel{15.} (p. 21.) Bulletin de la Société Française de Philosophie. IX. p. 274.

\notelabel{16.} (p. 21.) Se f. Ex. \textit{Matière et Mémoire}\textsuperscript{8}. p. 201 f. Ce qu’on appelle ordinairement un \textit{fait}, ce n’est pas la réalité telle qu’elle apparaîtrait à une \textit{intuition immédiate} . . . \textit{L’intuition pure, extérieure ou intérieure}, est celle d’une continuité indivisée. --- \textit{La Perception du Changement.} p. 8; 23. ---

\notelabel{17.} (p. 22.) Interessante Bilag til Belysning af Intuitionsbegrebet hos Bergson findes i to svenske Afhandlinger, den ene af Docent \emph{Malte Jacobsson} \textit{(Henri Bergsons Intuitionsfilosofi)}, den anden af Professor \emph{Hans Larsson} \textit{(Intuitionsproblemet.} Särskilt med hänsyn til Henri Bergsons Teori).

\notelabel{18.} (p. 24.) \emph{Descartes}: \textit{Règles pour la direction de l’esprit.} Chap. III et VII. (Oeuvres éd. Adam et Tannéry. Vol. X.). --- \emph{Henri Poincaré}: \textit{Science et Méthode.} Paris 1912. p. 47; 134; 157. --- Smlgn. ogsaa \emph{Kroman}’s „Anskuelsesskridt“ eller „Elementarvurderinger“ \textit{(Vor Naturerkendelse.} p. 97 f.). --- Det synes at have været Intuitioner af denne Art, der have ledet de ældste græske Matematikere til Opstilling af Sætninger, som først senere bleve Genstand for Bevisførelse. Se \emph{Zeuthen}: \textit{Den græske Geometri før Platon} (Videnskabernes Selskabs Oversigter 1914). --- Se \textit{Den nyere Filosofis Historie}\textsuperscript{2} I p. 211 f. \emph{Léon Brunschwieg}: \textit{Les Étapes de la Philosophie Mathématique.} Paris 1912. p. 105 ff.; 451. Brunschwieg betegner „den Modsætning, man i vore Dage har hævdet mellem Forstand og Intuition“ som „un accident malheureux de l’histoire“ (p. 452).% PRINTED AS IS: "Tannéry" (Tannery); "Brunschwieg" (Brunschvicg), three times

\notelabel{19.} (p. 25.) Smlgn. \textit{Den nyere Filosofis Historie}\textsuperscript{2} I p. 305 (og p. 328). --- \emph{Harold Joachim}: \textit{A Study of the Ethics of Spinoza.} Oxford 1901. p. 180--185. --- Det i Texten brugte Udtryk „Samfoldighed“ skyldes H. C. Ørsted. Se \textit{Danske Filosofer.} p. 51. --- Om den Intuition, der spiller en Rolle ved en Verdensanskuelses eller en Metafysiks Udvikling smlgn. \textit{Den menneskelige Tanke.} p. 306--315. --- Hvad \emph{Fichte} forstaar ved „intellektuel Anskuen“, er en umiddelbar Bevidsthed om vor egen Handlen, vor egen Selvvirksomhed. Enhver, der tilskriver sig selv en Virken, beraaber sig paa denne Anskuen. (Zweite Einleitung in die Wissenschaftslehre p. 5). (Smlgn. \textit{Den nyere Filosofis Historie}\textsuperscript{2} II p. 146). Den falder ikke sammen med Bergson’s Intuition, da den er en abstrakt Tanke, --- en Forudsætning, som kan \textit{uddrages} af umiddelbar Oplevelse, men
% --- p. 68 ---
\opage{68}som ikke falder helt sammen med nogen umiddelbar Oplevelse. Antagelsen af Aktivitet, af en Tilstands Udspring af vort eget Væsen og ikke af ydre Paavirkning, grunder sig paa \textit{Slutning}, ikke paa Anskuen, og Kant havde derfor med Rette nægtet, at der gaves „intellektuel Anskuen“.

\notelabel{20.} (p. 28.) \textit{Matière et Mémoire.}\textsuperscript{8} p. 251, smlgn. p. 4; 8; 66; 101.

\notelabel{21.} (p. 29.) \textit{Matière et Mémoire.}\textsuperscript{3} p. 65; 93; 149 f.

\notelabel{22.} (p. 29.) \textit{Matière et Mémoire.}\textsuperscript{3} p. 65; 251. --- Den Sprogbrug, ved hvilken de fysiologiske Processer, der foregaa i Sanseorganerne, betegnes som „perciperende“, forekommer ogsaa hos fysiologiske Forfattere, selv naar de (som \emph{Bechterew}: \textit{Objektive Psychologie.} Leipzig 1913 p. 68) udtrykkeligt erklære at ville se bort fra al subjektiv Iagttagelse. Man taler da f. Ex. om „perciperende Organer“, dels periferiske, dels centrale.

\notelabel{23.} (p. 30.) \textit{Matière et Mémoire.}\textsuperscript{3} p. 158. --- Med Hensyn til det Forudgaaende se p. 60; 159 (Artsforskel mellem Fornemmelse og Erindring).

\notelabel{24.} (p. 31.) \textit{Matière et Mémoire.}\textsuperscript{8} p. 244.

\notelabel{25.} (p. 31.) Se herom min \textit{Psykologi} V B, 4; D, 1--2. --- \textit{Den menneskelige Tanke} p. 99--108.

\notelabel{26.} (p. 32.) Spinozas Hypotese som Arbejdshypotese: Bulletin de la Société Française de Philosophie. I p. 68, --- som „constatation déguisée du fait“: Matière et Mémoire p. 247.

\notelabel{27.} (p. 34.) \emph{Bergson}: \textit{L’âme et le corps.} (I Skriftet Le Matérialisme actuel, der indeholder Foredrag af forskellige franske Forskere) p. 24 f.; 45--48. --- \emph{Kant’s} Udtalelse findes i \textit{Kritik der reinen Vernunft}\textsuperscript{2} p. 806--808. --- \emph{William James}: \textit{Human Immortality} (1899). --- Ytringen af \emph{Spinoza} findes \textit{Ethica} III, 2. Schol. ---

\notelabel{28.} (p. 35.) Smlgn. \textit{Den menneskelige Tanke.} p. 227--237.

\notelabel{29.} (p. 36.) \textit{Évolution créatrice.} Introduction.

\notelabel{30.} (p. 37.) \textit{Évolution créatrice.} p. 22; 101.

\notelabel{31.} (p. 39.) \textit{Évolution créatrice.} p. 49--50.

\notelabel{32.} (p. 40.) \textit{Évolution créatrice.} p. 137--152.

% --- p. 69 ---
\opage{69}\notelabel{33.} (p. 40.) \textit{Évolution créatrice.} p. 59--94.

\notelabel{34.} (p. 40.) \textit{Évolution créatrice.} p. 105--114; 273.

\notelabel{35.} (p. 41.) \emph{Jacques Loeb}: \textit{Das Leben.} 1911. p. 9.

\notelabel{36.} (p. 41.) \textit{Évolution créatrice} p. 194 (-rattacher % PRINTED AS IS: a hyphen after the parenthesis
la vie soit à la conscience, soit à quelque chose qui y ressemble); p. 279 (la vie est d’ordre psychologique). --- Smlgn. om Vitalisme min Psykologi II, 3 og \textit{Mindre Arbejder.} Første Række p. 40--50. --- Bergson har Forgængere i fransk Filosofi, især \emph{Félix Ravaisson}. Se dennes \textit{La philosophie en France au 19\textsuperscript{e} siècle} p. 246. (Om Spørgsmaalets tidligere Drøftelse i fransk Filosofi, se sammesteds p. 168--182).

\notelabel{37.} (p. 42.) Om Romantikens Udviklingsbegreb se \textit{Den nyere Filosofis Historie.}\textsuperscript{2} II p. 159--162. --- \textit{Mindre Arbejder.} Tredie Række. p. 205--207.

\notelabel{38.} (p. 44.) Smlgn. foruden min \textit{Psykologi} (Kap. IV og VII) mine \textit{Mindre Arbejder.} Tredie Række. p. 28--53. --- Bergson’s Stilling til Villiesproblemet har jeg allerede drøftet i \textit{Den menneskelige Tanke.} p. 291--295. --- Bergson’s Behandling af Problemet findes i \textit{Les Données immédiates de la Conscience.} Chap. 3. De l’organisation des états de conscience: la liberté. Allerede Titelen paa dette Afsnit er karakteristisk for hans Opfattelse.

\notelabel{39.} (p. 45.) Hvad Bergson kalder le moi fondamental eller le moi d’en bas, svarer ganske til, hvad jeg kalder „det reale Jeg“. Smlgn. \textit{Psykologi}\textsuperscript{6} p. 179--181; 354; 409--411.

\notelabel{40.} (p. 50.) Smlgn. min udførligere Fremstilling i min \textit{Psykologi} VI E, 9.

\notelabel{41.} (p. 53.) \textit{L’intuition philosophique.} (Revue de Métaphysique et de Morale. 1911. p. 813).

\notelabel{42.} (p. 54.) \textit{Bulletin de la Société Française de Philosophie.} VIII p. 341. \textit{Introduction à la Métaphysique.} (Revue de Métaphysique et de Morale. 1903). p. 25 (anført i „Den menneskelige Tanke“ p. 318). Flere Steder i \textit{Évolution créatrice} er Analogislutningen tydelig, se f. Ex. p. 190 f.; 250; 278--281; 387. --- \textit{La Perception du Changement} p. 7--9 bebrejder Bergson Filosoferne, at de bygge deres Systemer paa en enkelt Del eller Side af Tilværelsen, og han mener at undgaa denne Mislighed ved at fordybe sig i selve vor umiddelbare Opfattelse, den, af hvilken alle de stridende Begreber tilsidst udspringe.

\notelabel{43.} (p. 54.) Se nærmere om Analogier og deres Betydning \textit{Den menneskelige Tanke.} p. 181--186; 310--315.

% --- p. 70 ---
\opage{70}\notelabel{44.} (p. 55) \textit{Évolution créatrice.} p. 270; 323. --- Definitionen af Skabelse som Væxt findes p. 262.% PRINTED AS IS: the label reads "(p. 55)", without the stop

\notelabel{45.} (p. 57.) Se om Pragmatismen \textit{Den menneskelige Tanke.} p. 270--272.

\notelabel{46.} (p. 57.) \textit{Évolution créatrice.} p. 290 f. --- \textit{Introduction à la Métaphysique.} p. 19 f.

\notelabel{47.} (p. 58.) \textit{La Perception du Changement} p. 36.

\notelabel{48.} (p. 59.) Se nærmere om dette Punkt mine Helsingforsforelæsninger \textit{Personlighetsprincipen i Filosofin.} (Stockholm 1911) p. 7--38, og Afhandlingen \textit{Aandelig Kultur} (Mindre Arbejder. Tredie Række) p. 4--6.

\notelabel{49.} (p. 59.) \textit{Évolution créatrice} p. 280.

\notelabel{50.} (p. 59.) \textit{Évolution créatrice} p. 192. --- \textit{La Perception du Changement.} p. 9--13; 36. --- \textit{L’Intuition Philosophique} (Revue de Métaphysique et de Morale 1912). p. 827.

\notelabel{51.} (p. 60.) Nærmere herom i min \textit{Religionsfilosofi} \S{} 44; 115--116; 131.

\notelabel{52.} (p. 61.) Jeg kender kun denne Diskussion mellem Bergson og de jesuitiske Teologer gennem \emph{Marcel Hébert’s} Afhandling: „Henri Bergson et son affirmation de l’existence de Dieu“. (Coenobium. Rivista internationale di liberi studi. 1912).

\end{small}


\end{document}
